
\documentclass{article}
\usepackage{colortbl}
\usepackage{makecell}
\usepackage{multirow}
\usepackage{supertabular}

\begin{document}

\newcounter{utterance}

\centering \large Interaction Transcript for game `cryptolect', experiment `basic', episode 3 with qwen3.5{-}27b.
\vspace{24pt}

{ \footnotesize  \setcounter{utterance}{1}
\setlength{\tabcolsep}{0pt}
\begin{supertabular}{c@{$\;$}|p{.15\linewidth}@{}p{.15\linewidth}p{.15\linewidth}p{.15\linewidth}p{.15\linewidth}p{.15\linewidth}}
    \# & \multicolumn{2}{c}{Player} && \multicolumn{2}{c}{Game Master} \\
    \hline

    \theutterance \stepcounter{utterance}  
    & & & \multicolumn{4}{p{0.6\linewidth}}{
        \cellcolor[rgb]{0.9,0.9,0.9}{
            \makecell[{{p{\linewidth}}}]{
                \texttt{\tiny{[P1$\langle$GM]}}
                \texttt{You are a field linguist. You have made contact with the only speaker of an unknown} \\
\texttt{language (the "informant"). Nobody has ever written down this language before. Your task} \\
\texttt{is to work out how it works and then translate sentences into it correctly.} \\
\\ 
\texttt{You will study the language by asking the informant to translate English sentences for} \\
\texttt{you. You may only talk about a small fixed world, using exactly these English words:} \\
\\ 
\texttt{Nouns: boat, ship, hill, tree, stone, dog, farmer, road} \\
\texttt{Adjectives: (this language has no adjectives)} \\
\texttt{Verbs (write them in this base form): watch, hold, carry, hear, chase} \\
\texttt{Articles: the (definite), a (indefinite)} \\
\texttt{Number words: two, three (no number word = one / singular)} \\
\texttt{These are ALL the words in play {-} the language has no other words.} \\
\\ 
\texttt{HOW TO WRITE AN ENGLISH SENTENCE} \\
\texttt{Every sentence is one clause with a subject, a verb, and an object, written in this fixed English notation:} \\
\texttt{    SUBJECT  verb  OBJECT} \\
\texttt{Each of SUBJECT and OBJECT is a noun phrase of the form:} \\
\texttt{    {[}the | a{]}  {[}number word{]}  {[}adjective ...{]}  noun} \\
\texttt{{-} 'the' = definite, 'a' = indefinite (singular only); leave the article off for an indefinite plural.} \\
\texttt{{-} a number word (two, three) sets the count; with no number word the noun is singular.} \\
\texttt{Examples of well{-}formed English sentences you may send:} \\
\texttt{    a boat watch a ship} \\
\texttt{    three boats watch a ship} \\
\\ 
\texttt{WHAT YOU ALREADY KNOW} \\
\texttt{The informant has already translated these sentences for you. Each line is} \\
\texttt{"<English>  =>  <the language>":} \\
\texttt{    three trees watch two boats  =>  dugunu nu deniru gutinu ga} \\
\texttt{    three dogs hear a dog  =>  tipi puga tipinu ga} \\
\texttt{    three dogs carry three ships  =>  pirenu ga geti tipinu ga} \\
\texttt{    two dogs chase two boats  =>  dugunu nu buro tipinu nu} \\
\texttt{    two roads hold three trees  =>  gutinu ga nira dobudinu nu} \\
\texttt{    three farmers watch three farmers  =>  nidinu ga deniru nidinu ga} \\
\texttt{    three stones carry two roads  =>  dobudinu nu geti bodenu ga} \\
\texttt{    two trees hold three hills  =>  ruditinu ga nira gutinu nu} \\
\\ 
\texttt{HOW THE SESSION WORKS} \\
\texttt{1. INVESTIGATION: You have 4 rounds of questioning. In each round you may ask} \\
\texttt{   the informant to translate up to 3 English sentences (so at most} \\
\texttt{   12 probes in total). The informant answers each one truthfully and} \\
\texttt{   exactly. Use these answers to figure out the words and the grammar (word order, how} \\
\texttt{   number/case/definiteness are marked, how words must agree, and any sound changes).} \\
\texttt{   Choose your probes deliberately {-} test specific hypotheses.} \\
\texttt{2. EXAM: After the investigation rounds, the informant will give you 10 NEW} \\
\texttt{   English sentences. You must translate each one INTO the language. Every answer is} \\
\texttt{   graded by EXACT match against how the informant would say it {-} every word, ending and} \\
\texttt{   space must be right (capitalisation and surrounding spaces are ignored). The exam} \\
\texttt{   sentences will use the same vocabulary but in combinations you have not been shown, so} \\
\texttt{   you must induce the general rules, not just memorise examples.} \\
\\ 
\texttt{RESPONSE FORMAT (for every investigation round)} \\
\texttt{List your probe sentences under a 'PROBES:' line, one per line. A THOUGHTS line is} \\
\texttt{optional {-} keep any reasoning brief so you always finish writing your PROBES:} \\
\texttt{THOUGHTS: <optional; brief>} \\
\texttt{PROBES:} \\
\texttt{1. <an English sentence, written as described above>} \\
\texttt{2. <another English sentence>} \\
\texttt{(up to 3 lines; at least one)} \\
\\ 
\texttt{Begin your first round of probes now.} \\
            }
        }
    }
     \\ \\

    \theutterance \stepcounter{utterance}  
    & \multicolumn{4}{p{0.6\linewidth}}{
        \cellcolor[rgb]{0.561,0.737,0.561}{
            \makecell[{{p{\linewidth}}}]{
                \texttt{\tiny{[P1$\rangle$GM]}}
                \texttt{THOUGHTS: I need to isolate the vocabulary for the remaining nouns (hill, stone, farmer, road, dog, boat, ship, tree) and verbs (watch, hold, carry, hear, chase) that haven't been fully mapped or confirmed. I also need to determine the word order (SVO, SOV, VSO?), how definiteness ("the" vs "a") is marked, and how number (singular vs plural) and case (subject vs object) are handled.} \\
\\ 
\texttt{From the existing data:} \\
\texttt{{-} "three trees watch two boats" => "dugunu nu deniru gutinu ga"} \\
\texttt{{-} "three dogs hear a dog" => "tipi puga tipinu ga"} \\
\texttt{{-} "three dogs carry three ships" => "pirenu ga geti tipinu ga"} \\
\texttt{{-} "two dogs chase two boats" => "dugunu nu buro tipinu nu"} \\
\texttt{{-} "two roads hold three trees" => "gutinu ga nira dobudinu nu"} \\
\texttt{{-} "three farmers watch three farmers" => "nidinu ga deniru nidinu ga"} \\
\texttt{{-} "three stones carry two roads" => "dobudinu nu geti bodenu ga"} \\
\texttt{{-} "two trees hold three hills" => "ruditinu ga nira gutinu nu"} \\
\\ 
\texttt{Observations so far:} \\
\texttt{{-} "three" seems to be associated with "nu" or "ga" or part of the noun stem?} \\
\texttt{  {-} "three trees" = "dugunu" (maybe "dugu" + "nu"?)} \\
\texttt{  {-} "two boats" = "gutinu" (maybe "guti" + "nu"?)} \\
\texttt{  {-} "three dogs" = "tipi" (in "tipi puga tipinu ga") vs "tipinu" (in "pirenu ga geti tipinu ga"). Wait, "three dogs" is "tipi" in the subject position? Or is "tipi" the noun and "puga" the verb?} \\
\texttt{  {-} Let's look at "three dogs hear a dog": "tipi puga tipinu ga".} \\
\texttt{    {-} If "tipi" = "three dogs", then "puga" = "hear", "tipinu" = "a dog"?} \\
\texttt{    {-} But "three dogs carry three ships": "pirenu ga geti tipinu ga". Here "tipinu" is the object? Or subject?} \\
\texttt{    {-} "two dogs chase two boats": "dugunu nu buro tipinu nu".} \\
\texttt{      {-} "dugunu" = "two dogs"? But earlier "three trees" was "dugunu". Contradiction?} \\
\texttt{      {-} Let's re{-}examine "three trees watch two boats" => "dugunu nu deniru gutinu ga".} \\
\texttt{        {-} Maybe "dugunu" = "three trees"?} \\
\texttt{        {-} "gutinu" = "two boats"?} \\
\texttt{      {-} "two dogs chase two boats" => "dugunu nu buro tipinu nu".} \\
\texttt{        {-} If "dugunu" = "two dogs", then "dugunu" cannot be "three trees".} \\
\texttt{        {-} Unless "dugunu" is a number marker? No, "three trees" and "two dogs" are different.} \\
\texttt{        {-} Maybe the noun stem changes? Or the number marker is attached differently?} \\
\texttt{        {-} Let's look at "three farmers watch three farmers" => "nidinu ga deniru nidinu ga".} \\
\texttt{          {-} "nidinu" appears twice. "three farmers" = "nidinu"?} \\
\texttt{          {-} "deniru" = "watch"?} \\
\texttt{        {-} "three trees watch two boats" => "dugunu nu deniru gutinu ga".} \\
\texttt{          {-} "deniru" = "watch".} \\
\texttt{          {-} "dugunu" = "three trees"?} \\
\texttt{          {-} "gutinu" = "two boats"?} \\
\texttt{        {-} "two dogs chase two boats" => "dugunu nu buro tipinu nu".} \\
\texttt{          {-} "dugunu" = "two dogs"? This conflicts with "three trees" = "dugunu".} \\
\texttt{          {-} Unless "dugunu" is "two" + "something"?} \\
\texttt{          {-} Let's check "two roads hold three trees" => "gutinu ga nira dobudinu nu".} \\
\texttt{            {-} "gutinu" = "two roads"? But earlier "gutinu" was "two boats"?} \\
\texttt{            {-} "dobudinu" = "three trees"? But earlier "dugunu" was "three trees"?} \\
\texttt{          {-} This suggests the noun stem is changing based on number or case, or I am misidentifying the boundaries.} \\
\texttt{          {-} Let's try to find a consistent pattern.} \\
\texttt{          {-} "three trees" = "dugunu" (in sentence 1) vs "dobudinu" (in sentence 5)? No, sentence 5 is "two roads hold three trees" => "gutinu ga nira dobudinu nu".} \\
\texttt{            {-} If "dobudinu" = "three trees", then "dugunu" in sentence 1 must be something else.} \\
\texttt{            {-} Sentence 1: "three trees watch two boats" => "dugunu nu deniru gutinu ga".} \\
\texttt{            {-} Sentence 5: "two roads hold three trees" => "gutinu ga nira dobudinu nu".} \\
\texttt{            {-} Maybe "dugunu" = "three trees" and "dobudinu" = "three trees"? That's unlikely unless there's a typo in my reading or the language has variation.} \\
\texttt{            {-} Wait, "three trees" in sentence 1 is subject. "three trees" in sentence 5 is object.} \\
\texttt{            {-} Maybe the form changes based on case?} \\
\texttt{            {-} Let's look at "three dogs hear a dog" => "tipi puga tipinu ga".} \\
\texttt{              {-} "tipi" = "three dogs" (subject)?} \\
\texttt{              {-} "tipinu" = "a dog" (object)?} \\
\texttt{            {-} "three dogs carry three ships" => "pirenu ga geti tipinu ga".} \\
\texttt{              {-} "pirenu" = "three dogs"? But earlier "tipi" was "three dogs".} \\
\texttt{              {-} "tipinu" = "three ships"?} \\
\texttt{            {-} This is confusing. Let's look for a number marker.} \\
\texttt{            {-} "two" appears in: "two boats" (gutinu?), "two dogs" (dugunu?), "two roads" (gutinu?), "two roads" (dobudinu? no).} \\
\texttt{            {-} Let's assume the structure is {[}Subject{]} {[}Verb{]} {[}Object{]}.} \\
\texttt{            {-} Sentence 1: "three trees" (S) "watch" (V) "two boats" (O).} \\
\texttt{              {-} Words: dugunu, nu, deniru, gutinu, ga.} \\
\texttt{              {-} Maybe "dugunu" = "three trees", "deniru" = "watch", "gutinu" = "two boats".} \\
\texttt{              {-} What are "nu" and "ga"? Case markers?} \\
\texttt{            {-} Sentence 2: "three dogs" (S) "hear" (V) "a dog" (O).} \\
\texttt{              {-} Words: tipi, puga, tipinu, ga.} \\
\texttt{              {-} "tipi" = "three dogs"? "puga" = "hear"? "tipinu" = "a dog"? "ga" = case marker?} \\
\texttt{            {-} Sentence 3: "three dogs" (S) "carry" (V) "three ships" (O).} \\
\texttt{              {-} Words: pirenu, ga, geti, tipinu, ga.} \\
\texttt{              {-} "pirenu" = "three dogs"? "geti" = "carry"? "tipinu" = "three ships"?} \\
\texttt{              {-} But "tipinu" was "a dog" in sentence 2?} \\
\texttt{              {-} Maybe "tipi" is the stem for "dog", and "tipinu" is "dog" + something.} \\
\texttt{              {-} "tipi" (3 dogs) vs "tipinu" (1 dog)? No, "tipinu" is object in sentence 2 ("a dog") and object in sentence 3 ("three ships")? That doesn't make sense.} \\
\texttt{              {-} Let's re{-}read sentence 3: "three dogs carry three ships" => "pirenu ga geti tipinu ga".} \\
\texttt{                {-} Maybe "pirenu" = "three dogs", "geti" = "carry", "tipinu" = "three ships".} \\
\texttt{                {-} But "tipinu" was "a dog" in sentence 2?} \\
\texttt{                {-} Unless "tipi" is "dog" and "tipinu" is "ship"? No, "tipi" and "tipinu" look related.} \\
\texttt{                {-} Maybe "tipi" = "dog", "tipinu" = "ship"?} \\
\texttt{                {-} Sentence 2: "three dogs hear a dog" => "tipi puga tipinu ga".} \\
\texttt{                  {-} If "tipi" = "three dogs", "puga" = "hear", "tipinu" = "a dog".} \\
\texttt{                  {-} Then "tipinu" = "dog" (singular)?} \\
\texttt{                {-} Sentence 3: "three dogs carry three ships" => "pirenu ga geti tipinu ga".} \\
\texttt{                  {-} If "pirenu" = "three dogs", "geti" = "carry", "tipinu" = "three ships".} \\
\texttt{                  {-} Then "tipinu" = "ship"?} \\
\texttt{                  {-} Contradiction: "tipinu" cannot be both "dog" and "ship".} \\
\texttt{                  {-} Maybe "tipi" = "dog", "tipinu" = "dog" (plural/object)?} \\
\texttt{                  {-} Let's look at sentence 4: "two dogs chase two boats" => "dugunu nu buro tipinu nu".} \\
\texttt{                    {-} "dugunu" = "two dogs"?} \\
\texttt{                    {-} "buro" = "chase"?} \\
\texttt{                    {-} "tipinu" = "two boats"?} \\
\texttt{                    {-} But "tipinu" was "a dog" in sentence 2?} \\
\texttt{                    {-} This suggests "tipinu" is a noun stem that changes meaning? Unlikely.} \\
\texttt{                    {-} Maybe I am mis{-}segmenting.} \\
\texttt{                    {-} Let's try: "tipi" = "dog", "puga" = "hear", "tipinu" = "dog" (object form?).} \\
\texttt{                    {-} "pirenu" = "three dogs"? "geti" = "carry", "tipinu" = "three ships"?} \\
\texttt{                    {-} Maybe "tipinu" is "ship"? And "tipi" is "dog"?} \\
\texttt{                    {-} Sentence 2: "three dogs hear a dog" => "tipi puga tipinu ga".} \\
\texttt{                      {-} If "tipi" = "three dogs", "puga" = "hear", "tipinu" = "a dog".} \\
\texttt{                      {-} Then "tipinu" = "dog".} \\
\texttt{                    {-} Sentence 3: "three dogs carry three ships" => "pirenu ga geti tipinu ga".} \\
\texttt{                      {-} If "pirenu" = "three dogs", "geti" = "carry", "tipinu" = "three ships".} \\
\texttt{                      {-} Then "tipinu" = "ship".} \\
\texttt{                    {-} This is a contradiction. "tipinu" cannot be both "dog" and "ship".} \\
\texttt{                    {-} Unless "tipi" and "tipinu" are different words.} \\
\texttt{                    {-} Maybe "tipi" = "dog", "tipinu" = "ship"?} \\
\texttt{                    {-} Sentence 2: "three dogs hear a dog" => "tipi puga tipinu ga".} \\
\texttt{                      {-} "tipi" = "three dogs"? "puga" = "hear"? "tipinu" = "a dog"?} \\
\texttt{                      {-} If "tipinu" = "dog", then "tipi" = "three dogs"?} \\
\texttt{                      {-} But "tipi" and "tipinu" are similar. Maybe "tipi" is "dog" (singular) and "tipinu" is "dog" (plural)?} \\
\texttt{                      {-} Sentence 2: "three dogs" (plural) "hear" "a dog" (singular).} \\
\texttt{                        {-} "tipi" = "three dogs"? "tipinu" = "a dog"?} \\
\texttt{                        {-} If "tipi" is plural and "tipinu" is singular, that's odd.} \\
\texttt{                      {-} Let's look at sentence 4: "two dogs chase two boats" => "dugunu nu buro tipinu nu".} \\
\texttt{                        {-} "dugunu" = "two dogs"? "buro" = "chase"? "tipinu" = "two boats"?} \\
\texttt{                        {-} If "tipinu" = "boats", then in sentence 2 "tipinu" = "dog"?} \\
\texttt{                        {-} This implies "tipinu" is a generic noun? No.} \\
\texttt{                        {-} Maybe "tipi" = "dog", "tipinu" = "boat"?} \\
\texttt{                        {-} Sentence 2: "three dogs hear a dog" => "tipi puga tipinu ga".} \\
\texttt{                          {-} "tipi" = "three dogs"? "puga" = "hear"? "tipinu" = "a dog"?} \\
\texttt{                          {-} If "tipinu" = "dog", then "tipi" = "three dogs"?} \\
\texttt{                          {-} But "tipi" and "tipinu" are too similar.} \\
\texttt{                          {-} Maybe "tipi" = "dog", "tipinu" = "dog" (object)?} \\
\texttt{                          {-} Sentence 3: "three dogs carry three ships" => "pirenu ga geti tipinu ga".} \\
\texttt{                            {-} "pirenu" = "three dogs"? "geti" = "carry"? "tipinu" = "three ships"?} \\
\texttt{                            {-} If "tipinu" = "ships", then "tipi" = "dogs"?} \\
\texttt{                            {-} This suggests "tipi" = "dog", "tipinu" = "ship"?} \\
\texttt{                            {-} But sentence 2: "three dogs hear a dog" => "tipi puga tipinu ga".} \\
\texttt{                              {-} If "tipi" = "three dogs", "puga" = "hear", "tipinu" = "a dog".} \\
\texttt{                              {-} Then "tipinu" = "dog".} \\
\texttt{                              {-} Contradiction with sentence 3 where "tipinu" = "ships".} \\
\texttt{                              {-} Unless "tipinu" is "ship" and "tipi" is "dog", and sentence 2 is "three dogs hear a ship"? No, the English says "a dog".} \\
\texttt{                              {-} Maybe the translation is wrong? No, the informant is truthful.} \\
\texttt{                              {-} Maybe "tipi" = "dog", "tipinu" = "dog" (plural)?} \\
\texttt{                              {-} Sentence 2: "three dogs" (plural) "hear" "a dog" (singular).} \\
\texttt{                                {-} "tipi" = "three dogs"? "tipinu" = "a dog"?} \\
\texttt{                                {-} If "tipi" is plural and "tipinu" is singular, that's weird.} \\
\texttt{                              {-} Let's try a different approach. Maybe the number is marked on the verb or with a separate word.} \\
\texttt{                              {-} "three trees watch two boats" => "dugunu nu deniru gutinu ga".} \\
\texttt{                                {-} "dugunu" = "three trees"? "nu" = marker? "deniru" = "watch"? "gutinu" = "two boats"? "ga" = marker?} \\
\texttt{                              {-} "three dogs hear a dog" => "tipi puga tipinu ga".} \\
\texttt{                                {-} "tipi" = "three dogs"? "puga" = "hear"? "tipinu" = "a dog"? "ga" = marker?} \\
\texttt{                              {-} "three dogs carry three ships" => "pirenu ga geti tipinu ga".} \\
\texttt{                                {-} "pirenu" = "three dogs"? "ga" = marker? "geti" = "carry"? "tipinu" = "three ships"? "ga" = marker?} \\
\texttt{                              {-} "two dogs chase two boats" => "dugunu nu buro tipinu nu".} \\
\texttt{                                {-} "dugunu" = "two dogs"? "nu" = marker? "buro" = "chase"? "tipinu" = "two boats"? "nu" = marker?} \\
\texttt{                              {-} "two roads hold three trees" => "gutinu ga nira dobudinu nu".} \\
\texttt{                                {-} "gutinu" = "two roads"? "ga" = marker? "nira" = "hold"? "dobudinu" = "three trees"? "nu" = marker?} \\
\texttt{                              {-} "three farmers watch three farmers" => "nidinu ga deniru nidinu ga".} \\
\texttt{                                {-} "nidinu" = "three farmers"? "ga" = marker? "deniru" = "watch"? "nidinu" = "three farmers"? "ga" = marker?} \\
\texttt{                              {-} "three stones carry two roads" => "dobudinu nu geti bodenu ga".} \\
\texttt{                                {-} "dobudinu" = "three stones"? "nu" = marker? "geti" = "carry"? "bodenu" = "two roads"? "ga" = marker?} \\
\texttt{                              {-} "two trees hold three hills" => "ruditinu ga nira gutinu nu".} \\
\texttt{                                {-} "ruditinu" = "two trees"? "ga" = marker? "nira" = "hold"? "gutinu" = "three hills"? "nu" = marker?} \\
\\ 
\texttt{Let's look for patterns in the endings:} \\
\texttt{{-} "nu" and "ga" seem to be case markers or number markers.} \\
\texttt{{-} In sentence 1: "dugunu nu" (S) ... "gutinu ga" (O).} \\
\texttt{{-} In sentence 2: "tipi" (S) ... "tipinu ga" (O).} \\
\texttt{{-} In sentence 3: "pirenu ga" (S) ... "tipinu ga" (O).} \\
\texttt{{-} In sentence 4: "dugunu nu" (S) ... "tipinu nu" (O).} \\
\texttt{{-} In sentence 5: "gutinu ga" (S) ... "dobudinu nu" (O).} \\
\texttt{{-} In sentence 6: "nidinu ga" (S) ... "nidinu ga" (O).} \\
\texttt{{-} In sentence 7: "dobudinu nu" (S) ... "bodenu ga" (O).} \\
\texttt{{-} In sentence 8: "ruditinu ga" (S) ... "gutinu nu" (O).} \\
\\ 
\texttt{Hypothesis:} \\
\texttt{{-} "nu" and "ga" are case markers.} \\
\texttt{{-} Subject case: sometimes "nu", sometimes "ga".} \\
\texttt{{-} Object case: sometimes "ga", sometimes "nu".} \\
\texttt{{-} Let's check the number:} \\
\texttt{  {-} "three trees" = "dugunu" (S) in sentence 1.} \\
\texttt{  {-} "two boats" = "gutinu" (O) in sentence 1.} \\
\texttt{  {-} "three dogs" = "tipi" (S) in sentence 2.} \\
\texttt{  {-} "a dog" = "tipinu" (O) in sentence 2.} \\
\texttt{  {-} "three dogs" = "pirenu" (S) in sentence 3.} \\
\texttt{  {-} "three ships" = "tipinu" (O) in sentence 3.} \\
\texttt{  {-} "two dogs" = "dugunu" (S) in sentence 4.} \\
\texttt{  {-} "two boats" = "tipinu" (O) in sentence 4.} \\
\texttt{  {-} "two roads" = "gutinu" (S) in sentence 5.} \\
\texttt{  {-} "three trees" = "dobudinu" (O) in sentence 5.} \\
\texttt{  {-} "three farmers" = "nidinu" (S) in sentence 6.} \\
\texttt{  {-} "three farmers" = "nidinu" (O) in sentence 6.} \\
\texttt{  {-} "three stones" = "dobudinu" (S) in sentence 7.} \\
\texttt{  {-} "two roads" = "bodenu" (O) in sentence 7.} \\
\texttt{  {-} "two trees" = "ruditinu" (S) in sentence 8.} \\
\texttt{  {-} "three hills" = "gutinu" (O) in sentence 8.} \\
\\ 
\texttt{This is very confusing. The same English phrase maps to different words:} \\
\texttt{{-} "three trees": "dugunu" (S1), "dobudinu" (O5).} \\
\texttt{{-} "two boats": "gutinu" (O1), "tipinu" (O4).} \\
\texttt{{-} "three dogs": "tipi" (S2), "pirenu" (S3).} \\
\texttt{{-} "two dogs": "dugunu" (S4).} \\
\texttt{{-} "two roads": "gutinu" (S5), "bodenu" (O7).} \\
\texttt{{-} "three trees": "dobudinu" (O5).} \\
\texttt{{-} "three hills": "gutinu" (O8).} \\
\\ 
\texttt{Wait, maybe the noun stem changes based on number?} \\
\texttt{{-} "tree": "dugu" (3), "dobudu" (3)? No.} \\
\texttt{{-} "boat": "guti" (2), "tipi" (2)? No.} \\
\texttt{{-} "dog": "tipi" (3), "pire" (3), "dugu" (2)? No.} \\
\\ 
\texttt{Let's look at the endings again.} \\
\texttt{{-} "dugunu", "gutinu", "tipinu", "pirenu", "dobudinu", "nidinu", "bodenu", "ruditinu".} \\
\texttt{{-} All end in "nu" or "inu" or "enu".} \\
\texttt{{-} "tipi" (no ending) in sentence 2.} \\
\texttt{{-} "pirenu" in sentence 3.} \\
\texttt{{-} "dugunu" in sentence 1 and 4.} \\
\texttt{{-} "gutinu" in sentence 1, 5, 8.} \\
\texttt{{-} "tipinu" in sentence 2, 3, 4.} \\
\texttt{{-} "dobudinu" in sentence 5, 7.} \\
\texttt{{-} "nidinu" in sentence 6.} \\
\texttt{{-} "bodenu" in sentence 7.} \\
\texttt{{-} "ruditinu" in sentence 8.} \\
\\ 
\texttt{Maybe the number is part of the word?} \\
\texttt{{-} "three" = "nu"? "two" = "nu"? No, both have "nu".} \\
\texttt{{-} "three" = "dugu", "two" = "guti"?} \\
\texttt{{-} Let's assume the structure is {[}Noun Stem + Number + Case{]}.} \\
\texttt{{-} "three trees" = "dugunu" (S1). "dugu" = tree, "nu" = three?} \\
\texttt{{-} "two boats" = "gutinu" (O1). "guti" = boat, "nu" = two?} \\
\texttt{{-} "three dogs" = "tipi" (S2). "tipi" = dog, no number marker? But it's "three".} \\
\texttt{{-} "a dog" = "tipinu" (O2). "tipi" = dog, "nu" = one?} \\
\texttt{{-} "three dogs" = "pirenu" (S3). "pire" = dog, "nu" = three?} \\
\texttt{{-} "three ships" = "tipinu" (O3). "tipi" = ship, "nu" = three?} \\
\texttt{{-} "two dogs" = "dugunu" (S4). "dugu" = dog, "nu} \\
            }
        }
    }
    & & \\ \\

    \theutterance \stepcounter{utterance}  
    & & & \multicolumn{2}{p{0.3\linewidth}}{
        \cellcolor[rgb]{0.9,0.9,0.9}{
            \makecell[{{p{\linewidth}}}]{
                \texttt{\tiny{[GM$|$GM]}}
                \texttt{could not find any probe sentences; put them under a 'PROBES:' line, one per line} \\
            }
        }
    }
    & & \\ \\

    \theutterance \stepcounter{utterance}  
    & & & \multicolumn{4}{p{0.6\linewidth}}{
        \cellcolor[rgb]{0.9,0.9,0.9}{
            \makecell[{{p{\linewidth}}}]{
                \texttt{\tiny{[P1$\langle$GM]}}
                \texttt{Your last response could not be processed: could not find any probe sentences; put them under a 'PROBES:' line, one per line.} \\
\texttt{List your probes (a brief THOUGHTS line is optional):} \\
\texttt{PROBES:} \\
\texttt{1. <an English sentence>} \\
\texttt{(up to 3 probe lines; at least one)} \\
            }
        }
    }
     \\ \\

    \theutterance \stepcounter{utterance}  
    & \multicolumn{4}{p{0.6\linewidth}}{
        \cellcolor[rgb]{0.561,0.737,0.561}{
            \makecell[{{p{\linewidth}}}]{
                \texttt{\tiny{[P1$\rangle$GM]}}
                \texttt{THOUGHTS: The previous data is confusing because the same English phrases map to different words (e.g., "three trees" is "dugunu" in one sentence and "dobudinu" in another). This suggests the noun stem changes based on number, case, or definiteness, or I am mis{-}segmenting. I need to isolate singular forms to establish base stems. I will probe with singular subjects and objects using "a" (indefinite) and "the" (definite) to see how definiteness and singular number are marked, and to get clean base forms for "dog", "boat", "ship", "tree", "hill", "stone", "road", "farmer". I will also test the verb "watch" with singulars to confirm the verb form.} \\
\\ 
\texttt{PROBES:} \\
\texttt{1. a dog watch a boat} \\
\texttt{2. the dog watch the boat} \\
\texttt{3. a tree hold a hill} \\
            }
        }
    }
    & & \\ \\

    \theutterance \stepcounter{utterance}  
    & & & \multicolumn{2}{p{0.3\linewidth}}{
        \cellcolor[rgb]{0.9,0.9,0.9}{
            \makecell[{{p{\linewidth}}}]{
                \texttt{\tiny{[GM$|$GM]}}
                {'kind': 'investigation', 'round': 0, 'probes_issued': 3, 'malformed': 0, 'probes': ['a dog watch a boat', 'the dog watch the boat', 'a tree hold a hill']}
            }
        }
    }
    & & \\ \\

    \theutterance \stepcounter{utterance}  
    & & & \multicolumn{4}{p{0.6\linewidth}}{
        \cellcolor[rgb]{0.9,0.9,0.9}{
            \makecell[{{p{\linewidth}}}]{
                \texttt{\tiny{[P1$\langle$GM]}}
                \texttt{TRANSLATIONS:} \\
\texttt{1. "a dog watch a boat"  =>  dugu deniru tipi} \\
\texttt{2. "the dog watch the boat"  =>  dugu deniru tipi} \\
\texttt{3. "a tree hold a hill"  =>  ruditi nira guti} \\
\texttt{STATUS: investigation round 1 of 4 done. Probes used: 3 of 12. 3 investigation round(s) left before the exam (10 sentences).} \\
\texttt{Send your next probes now (a PROBES list; brief THOUGHTS optional).} \\
            }
        }
    }
     \\ \\

    \theutterance \stepcounter{utterance}  
    & \multicolumn{4}{p{0.6\linewidth}}{
        \cellcolor[rgb]{0.561,0.737,0.561}{
            \makecell[{{p{\linewidth}}}]{
                \texttt{\tiny{[P1$\rangle$GM]}}
                \texttt{THOUGHTS: The new data is crucial.} \\
\texttt{1. "a dog watch a boat" => "dugu deniru tipi"} \\
\texttt{2. "the dog watch the boat" => "dugu deniru tipi"} \\
\texttt{   {-} Conclusion: Definiteness ("the" vs "a") is **not** marked in the language. The form is identical.} \\
\texttt{   {-} "dugu" = dog (singular).} \\
\texttt{   {-} "deniru" = watch.} \\
\texttt{   {-} "tipi" = boat (singular).} \\
\texttt{   {-} Word order: Subject {-} Verb {-} Object (SVO).} \\
\texttt{   {-} No case markers for singulars? Or maybe the markers "nu"/"ga" are only for plurals?} \\
\\ 
\texttt{3. "a tree hold a hill" => "ruditi nira guti"} \\
\texttt{   {-} "ruditi" = tree (singular).} \\
\texttt{   {-} "nira" = hold.} \\
\texttt{   {-} "guti" = hill (singular).} \\
\\ 
\texttt{Now let's re{-}evaluate the plural data with these base stems:} \\
\texttt{{-} "three trees" (S) in "three trees watch two boats" => "dugunu nu deniru gutinu ga".} \\
\texttt{  {-} Base "tree" = "ruditi". Plural "three trees" = "dugunu"? That's a huge change.} \\
\texttt{  {-} Wait, "dugunu" in the new data is "dog" (singular).} \\
\texttt{  {-} In sentence 1: "three trees watch two boats" => "dugunu nu deniru gutinu ga".} \\
\texttt{    {-} If "dugunu" = "dog" (singular), then the translation "three trees watch two boats" => "dog ... watch ... boat" is wrong.} \\
\texttt{    {-} Unless "dugunu" is NOT "dog".} \\
\texttt{    {-} Let's re{-}read the new translation: "a dog watch a boat" => "dugu deniru tipi".} \\
\texttt{      {-} "dugu" = dog.} \\
\texttt{      {-} "tipi" = boat.} \\
\texttt{    {-} Old sentence 1: "three trees watch two boats" => "dugunu nu deniru gutinu ga".} \\
\texttt{      {-} "deniru" = watch.} \\
\texttt{      {-} "dugunu" must be "three trees".} \\
\texttt{      {-} "gutinu" must be "two boats".} \\
\texttt{    {-} But "dugu" is "dog". "dugunu" looks like "dugu" + "nu".} \\
\texttt{    {-} Is it possible "dugu" means "tree" and "dog" is something else?} \\
\texttt{    {-} Let's check "a dog watch a boat" => "dugu deniru tipi".} \\
\texttt{      {-} If "dugu" = tree, then "a tree watch a boat". But the English is "a dog".} \\
\texttt{      {-} So "dugu" = dog.} \\
\texttt{    {-} Then "dugunu" in sentence 1 must be "three trees".} \\
\texttt{    {-} This implies "dugunu" is not derived from "dugu" (dog).} \\
\texttt{    {-} Maybe "dugunu" = "three" + "tree"?} \\
\texttt{    {-} Let's look at "three trees" in sentence 5: "two roads hold three trees" => "gutinu ga nira dobudinu nu".} \\
\texttt{      {-} "dobudinu" = "three trees".} \\
\texttt{      {-} So "three trees" is "dugunu" in sentence 1 and "dobudinu" in sentence 5?} \\
\texttt{      {-} This suggests the noun stem changes based on case (Subject vs Object) or number?} \\
\texttt{      {-} Sentence 1: "three trees" (Subject) = "dugunu".} \\
\texttt{      {-} Sentence 5: "three trees" (Object) = "dobudinu".} \\
\texttt{      {-} This is a strong pattern: Subject form vs Object form.} \\
\texttt{      {-} Let's check "two boats":} \\
\texttt{        {-} Sentence 1: "two boats" (Object) = "gutinu".} \\
\texttt{        {-} Sentence 4: "two dogs chase two boats" => "dugunu nu buro tipinu nu".} \\
\texttt{          {-} "tipinu" = "two boats" (Object).} \\
\texttt{          {-} So "two boats" is "gutinu" in S1 and "tipinu" in S4?} \\
\texttt{          {-} Wait, "tipi" is "boat" (singular) from the new data.} \\
\texttt{          {-} "tipinu" = "tipi" + "nu".} \\
\texttt{          {-} "gutinu" = "guti" + "nu"? "guti" is "hill" (singular) from new data.} \\
\texttt{          {-} This is very confusing. "guti" = hill. "gutinu" = two boats?} \\
\texttt{          {-} Maybe "guti" is not "hill"?} \\
\texttt{          {-} "a tree hold a hill" => "ruditi nira guti".} \\
\texttt{            {-} "ruditi" = tree.} \\
\texttt{            {-} "nira" = hold.} \\
\texttt{            {-} "guti" = hill.} \\
\texttt{          {-} "two roads hold three trees" => "gutinu ga nira dobudinu nu".} \\
\texttt{            {-} "gutinu" = "two roads" (Subject).} \\
\texttt{            {-} "dobudinu" = "three trees" (Object).} \\
\texttt{          {-} "two trees hold three hills" => "ruditinu ga nira gutinu nu".} \\
\texttt{            {-} "ruditinu" = "two trees" (Subject).} \\
\texttt{            {-} "gutinu" = "three hills" (Object).} \\
\texttt{          {-} So "gutinu" is "two roads" (S) and "three hills" (O)?} \\
\texttt{          {-} This implies the word "gutinu" changes meaning based on context? No, that's impossible.} \\
\texttt{          {-} Maybe I have the base stems wrong.} \\
\texttt{          {-} Let's re{-}examine "a tree hold a hill" => "ruditi nira guti".} \\
\texttt{            {-} Maybe "ruditi" = hill? "guti" = tree?} \\
\texttt{            {-} "two trees hold three hills" => "ruditinu ga nira gutinu nu".} \\
\texttt{              {-} If "ruditinu" = "two trees", then "ruditi" = tree.} \\
\texttt{              {-} If "gutinu" = "three hills", then "guti" = hill.} \\
\texttt{            {-} So "ruditi" = tree, "guti" = hill.} \\
\texttt{          {-} Now "two roads hold three trees" => "gutinu ga nira dobudinu nu".} \\
\texttt{            {-} "gutinu" = "two roads".} \\
\texttt{            {-} "dobudinu" = "three trees".} \\
\texttt{            {-} But "gutinu" was "three hills" in the previous sentence?} \\
\texttt{            {-} "two trees hold three hills" => "ruditinu ga nira gutinu nu".} \\
\texttt{              {-} "gutinu" = "three hills".} \\
\texttt{            {-} "two roads hold three trees" => "gutinu ga nira dobudinu nu".} \\
\texttt{              {-} "gutinu" = "two roads".} \\
\texttt{            {-} This is a contradiction. "gutinu" cannot be both "three hills" and "two roads".} \\
\texttt{            {-} Unless "gutinu" is a number marker? No, it's attached to the noun.} \\
\texttt{            {-} Maybe the noun stem for "road" is "guti"?} \\
\texttt{            {-} "two roads" = "gutinu".} \\
\texttt{            {-} "three hills" = "gutinu"?} \\
\texttt{            {-} This implies "road" and "hill" have the same stem? Unlikely.} \\
\texttt{            {-} Let's look at the endings again.} \\
\texttt{            {-} "ruditinu" (2 trees S), "dobudinu" (3 trees O), "dugunu" (3 trees S), "gutinu" (2 boats O), "tipinu" (2 boats O), "gutinu" (2 roads S), "gutinu" (3 hills O).} \\
\texttt{            {-} There is a pattern of "nu" and "ga" as case markers.} \\
\texttt{            {-} S1: "dugunu nu" (S) ... "gutinu ga" (O).} \\
\texttt{            {-} S4: "dugunu nu" (S) ... "tipinu nu" (O).} \\
\texttt{            {-} S5: "gutinu ga" (S) ... "dobudinu nu" (O).} \\
\texttt{            {-} S8: "ruditinu ga" (S) ... "gutinu nu" (O).} \\
\texttt{            {-} It seems "nu" and "ga" are case markers, but their assignment to Subject/Object is not consistent.} \\
\texttt{            {-} Maybe "nu" is Subject and "ga" is Object?} \\
\texttt{              {-} S1: S="dugunu nu", O="gutinu ga". (Matches)} \\
\texttt{              {-} S4: S="dugunu nu", O="tipinu nu". (O has "nu", not "ga").} \\
\texttt{              {-} S5: S="gutinu ga", O="dobudinu nu". (S has "ga", O has "nu").} \\
\texttt{              {-} S8: S="ruditinu ga", O="gutinu nu". (S has "ga", O has "nu").} \\
\texttt{            {-} So "nu" and "ga" are not fixed to Subject/Object.} \\
\texttt{            {-} Maybe they depend on the number?} \\
\texttt{              {-} "three" vs "two"?} \\
\texttt{              {-} S1: 3 trees (S) {-}> "nu", 2 boats (O) {-}> "ga".} \\
\texttt{              {-} S4: 2 dogs (S) {-}> "nu", 2 boats (O) {-}> "nu".} \\
\texttt{              {-} S5: 2 roads (S) {-}> "ga", 3 trees (O) {-}> "nu".} \\
\texttt{              {-} S8: 2 trees (S) {-}> "ga", 3 hills (O) {-}> "nu".} \\
\texttt{            {-} This is messy.} \\
\texttt{            {-} Let's try to isolate the number words.} \\
\texttt{            {-} "three" might be "nu"? "two" might be "ga"?} \\
\texttt{            {-} S1: "three trees" = "dugunu". "two boats" = "gutinu".} \\
\texttt{            {-} S4: "two dogs" = "dugunu". "two boats" = "tipinu".} \\
\texttt{            {-} S5: "two roads" = "gutinu". "three trees" = "dobudinu".} \\
\texttt{            {-} S8: "two trees" = "ruditinu". "three hills" = "gutinu".} \\
\texttt{            {-} Notice "gutinu" appears as "two boats" (O), "two roads" (S), "three hills" (O).} \\
\texttt{            {-} This suggests "gutinu" is not a single noun+number.} \\
\texttt{            {-} Maybe the noun stem is "guti" and the number/case is "nu"?} \\
\texttt{            {-} "guti" = boat? road? hill?} \\
\texttt{            {-} "a tree hold a hill" => "ruditi nira guti".} \\
\texttt{              {-} "guti" = hill.} \\
\texttt{            {-} "two trees hold three hills" => "ruditinu ga nira gutinu nu".} \\
\texttt{              {-} "gutinu" = "three hills".} \\
\texttt{              {-} So "guti" + "nu" = "three hills"?} \\
\texttt{              {-} But "guti" is singular "hill". "gutinu" is "three hills".} \\
\texttt{              {-} So "nu" = "three"?} \\
\texttt{            {-} "two roads hold three trees" => "gutinu ga nira dobudinu nu".} \\
\texttt{              {-} "gutinu" = "two roads".} \\
\texttt{              {-} If "nu" = "three", then "gutinu" cannot be "two roads".} \\
\texttt{              {-} Unless "gutinu" = "two" + "road"?} \\
\texttt{              {-} But "gutinu" was "three hills" in S8.} \\
\texttt{              {-} This implies "guti" is not "hill".} \\
\texttt{              {-} Let's re{-}evaluate "a tree hold a hill" => "ruditi nira guti".} \\
\texttt{                {-} Maybe "ruditi" = hill, "guti" = tree?} \\
\texttt{                {-} "two trees hold three hills" => "ruditinu ga nira gutinu nu".} \\
\texttt{                  {-} "ruditinu" = "two trees".} \\
\texttt{                  {-} "gutinu" = "three hills".} \\
\texttt{                  {-} If "ruditinu" = "two trees", then "ruditi" = tree.} \\
\texttt{                  {-} If "gutinu" = "three hills", then "guti" = hill.} \\
\texttt{                {-} So "ruditi" = tree, "guti" = hill.} \\
\texttt{                {-} Now "two roads hold three trees" => "gutinu ga nira dobudinu nu".} \\
\texttt{                  {-} "gutinu" = "two roads".} \\
\texttt{                  {-} "dobudinu" = "three trees".} \\
\texttt{                  {-} But "gutinu" was "three hills" in S8.} \\
\texttt{                  {-} Contradiction: "gutinu" cannot be "two roads" and "three hills".} \\
\texttt{                  {-} Unless "road" and "hill" are the same word? No.} \\
\texttt{                  {-} Maybe the translation of "two roads hold three trees" is "gutinu ga nira dobudinu nu" where "gutinu" = "two roads" and "dobudinu" = "three trees".} \\
\texttt{                  {-} And "two trees hold three hills" is "ruditinu ga nira gutinu nu" where "ruditinu" = "two trees" and "gutinu" = "three hills".} \\
\texttt{                  {-} This means "gutinu" is used for "two roads" and "three hills".} \\
\texttt{                  {-} This is impossible unless the language has homonyms or I am missing a crucial rule.} \\
\texttt{                  {-} Wait, maybe "gutinu" is "two" + "something" and "gutinu" is "three" + "something"?} \\
\texttt{                  {-} No, the word is the same.} \\
\texttt{                  {-} Let's look at the endings again.} \\
\texttt{                  {-} "ruditinu" (2 trees S), "dobudinu" (3 trees O), "dugunu" (3 trees S), "gutinu" (2 boats O), "tipinu" (2 boats O), "gutinu" (2 roads S), "gutinu" (3 hills O).} \\
\texttt{                  {-} Maybe the noun stem is "guti" and the number is encoded in the case marker?} \\
\texttt{                  {-} Or maybe "gutinu" is "two" + "road" and "gutinu" is "three" + "hill"?} \\
\texttt{                  {-} This suggests the number is not a separate word but part of the noun form, and the form depends on the number.} \\
\texttt{                  {-} But "gutinu" is used for "two roads" and "three hills".} \\
\texttt{                  {-} This implies "road" and "hill" have the same stem "guti"?} \\
\texttt{                  {-} "a tree hold a hill" => "ruditi nira guti".} \\
\texttt{                  {-} "two roads hold three trees" => "gutinu ga nira dobudinu nu".} \\
\texttt{                  {-} If "guti" = hill, then "gutinu" = "three hills" (S8) and "two roads" (S5)?} \\
\texttt{                  {-} This is a contradiction.} \\
\texttt{                  {-} Let's try a different hypothesis: The word order is not SVO.} \\
\texttt{                  {-} Maybe it's SOV?} \\
\texttt{                  {-} "three trees watch two boats" => "dugunu nu deniru gutinu ga".} \\
\texttt{                    {-} If SOV: "dugunu nu" (S) "gutinu ga" (O) "deniru" (V).} \\
\texttt{                    {-} "deniru" = watch.} \\
\texttt{                  {-} "a dog watch a boat" => "dugu deniru tipi".} \\
\texttt{                    {-} If SOV: "dugu" (S) "tipi" (O) "deniru" (V).} \\
\texttt{                    {-} But the translation is "dugu deniru tipi".} \\
\texttt{                    {-} So it's SVO.} \\
\texttt{                  {-} Let's go back to the contradiction.} \\
\texttt{                  {-} "gutinu" = "two roads" (S5) and "three hills" (S8).} \\
\texttt{                  {-} Maybe "gutinu" is "two" + "road" and "gutinu" is "three" + "hill"?} \\
\texttt{                  {-} This implies the number is not marked on the noun, but the noun stem changes based on number?} \\
\texttt{                  {-} Or maybe "gutinu" is a number word?} \\
\texttt{                  {-} "gutinu" = "two"? "gutinu" = "three"?} \\
\texttt{                  {-} No, "gutinu" is attached to the noun.} \\
\texttt{                  {-} Let's look at "dugunu".} \\
\texttt{                  {-} "three trees" (S1) = "dugunu".} \\
\texttt{                  {-} "two dogs" (S4) = "dugunu".} \\
\texttt{                  {-} "dugunu" = "three trees" and "two dogs".} \\
\texttt{                  {-} This implies "dugunu" is not a noun+number.} \\
\texttt{                  {-} Maybe "dugunu" is "three" + "tree" and "two" + "dog"?} \\
\texttt{                  {-} This suggests the number is part of the word, and the word changes based on the number.} \\
\texttt{                  {-} But "dugunu" is used for "three trees" and "two dogs".} \\
\texttt{                  {-} This is impossible unless "tree" and "dog" have the same stem "dugu"?} \\
\texttt{                  {-} "a dog watch a boat" => "dugu deniru tipi".} \\
\texttt{                    {-} "dugu" = dog.} \\
\texttt{                  {-} "three trees watch two boats" => "dugunu nu deniru gutinu ga".} \\
\texttt{                    {-} "dugunu" = "three trees".} \\
\texttt{                    {-} If "dugu" = dog, then "dugunu" = "dog" + "nu".} \\
\texttt{                    {-} But "dugunu" = "three trees".} \\
\texttt{                    {-} This implies "tree" = "dog"? No.} \\
\texttt{                  {-} Maybe "dugunu" is "three" + "tree" and "dugu" is "dog".} \\
\texttt{                  {-} But "dugunu" is also "two dogs".} \\
\texttt{                  {-} This is a mess.} \\
\texttt{                  {-} Let's try to find a pattern in the endings.} \\
\texttt{                  {-} "nu" and "ga" are case markers.} \\
\texttt{                  {-} "dugunu" = "dugu" + "nu".} \\
\texttt{                  {-} "gutinu" = "guti" + "nu".} \\
\texttt{                  {-} "tipinu" = "tipi" + "nu".} \\
\texttt{                  {-} "dobudinu" = "dobudu" + "nu".} \\
\texttt{                  {-} "nidinu" = "nidi" + "nu".} \\
\texttt{                  {-} "bodenu" = "bode" + "nu".} \\
\texttt{                  {-} "ruditinu" = "ruditi" + "nu".} \\
\texttt{                  {-} "pirenu" = "pire" + "nu".} \\
\texttt{                  {-} "tipi" (no ending) in S2.} \\
\texttt{                  {-} "dugu" (no ending) in new data.} \\
\texttt{                  {-} "ruditi" (no ending) in new data.} \\
\texttt{                  {-} "guti" (no ending) in new data.} \\
\texttt{                  {-} So singular nouns have no ending. Plural nouns have "nu" or "inu" or "enu".} \\
\texttt{                  {-} "dugunu" = "dugu" + "nu".} \\
\texttt{                  {-} "gutinu" = "guti" + "nu".} \\
\texttt{                  {-} "tipinu" = "tipi" + "nu".} \\
\texttt{                  {-} "dobudinu" = "dobudu" + "nu".} \\
\texttt{                  {-} "nidinu" = "nidi" + "nu".} \\
\texttt{                  {-} "bodenu" = "bode" + "nu".} \\
\texttt{                  {-} "ruditinu" = "ruditi" + "nu".} \\
\texttt{                  {-} "pirenu" = "pire" + "nu".} \\
\texttt{                  {-} So the plural marker is "nu" (or "inu", "enu" depending on the stem).} \\
\texttt{                  {-} Now, what about the number "two" vs "three"?} \\
\texttt{                  {-} "three trees" = "dugunu" (S1).} \\
\texttt{                  {-} "two dogs" = "dugunu" (S4).} \\
\texttt{                  {-} "dugunu" is used for both "three trees" and "two dogs".} \\
\texttt{                  {-} This implies "dugunu" is not "three trees" or "two dogs".} \\
\texttt{                  {-} Maybe "dugunu" is "three" + "tree" and "two" + "dog"?} \\
\texttt{                  {-} But "dugunu" is the same word.} \\
\texttt{                  {-} This suggests the number is not marked on the noun, but the noun stem changes based on the number?} \\
\texttt{                  {-} Or maybe "dugunu" is "three" and "tree" is "dugu"?} \\
\texttt{                  {-} "three trees" = "dugunu".} \\
\texttt{                  {-} "two dogs" = "dugunu".} \\
\texttt{                  {-} This implies "tree" = "dog"? No.} \\
\texttt{                  {-} Maybe "dugunu" is "three" + "something" and "two" + "something"?} \\
\texttt{                  {-} Let's look at "three dogs" in S2: "tipi".} \\
\texttt{                  {-} "three dogs" in S3: "pirenu".} \\
\texttt{                  {-} "tipi" = "three dogs"? "pirenu" = "three dogs"?} \\
\texttt{                  {-} This is inconsistent.} \\
\texttt{                  {-} Maybe "tipi" is "dog" and "pirenu" is "three dogs"?} \\
\texttt{                  {-} But "tipi" is "three dogs" in S2.} \\
\texttt{                  {-} "three dogs hear a dog" => "tipi puga tipinu ga".} \\
\texttt{                    {-} "tipi" = "three dogs".} \\
\texttt{                    {-} "puga" = "hear".} \\
\texttt{                    {-} "tipinu" = "a dog".} \\
\texttt{                  {-} "three dogs carry three ships" => "pirenu ga geti tipinu ga".} \\
\texttt{                    {-} "pirenu" = "three dogs".} \\
\texttt{                    {-} "geti" = "carry".} \\
\texttt{                    {-} "tipinu" = "three ships".} \\
\texttt{                  {-} So "tipi" = "three dogs" and "pirenu" = "three dogs".} \\
\texttt{                  {-} This implies "tipi" and "pirenu" are synonyms? Or "tipi" is "dog" and "pirenu" is "three dogs"?} \\
\texttt{                  {-} But "tipi" is "three dogs" in S2.} \\
\texttt{                  {-} Maybe "tipi" is "dog" and the number is marked elsewhere?} \\
\texttt{                  {-} "three dogs" = "tipi" (S2).} \\
\texttt{                  {-} "three dogs" = "pirenu" (S3).} \\
\texttt{                  {-} This is a contradiction.} \\
\texttt{                  {-} Unless "tipi" is "dog" and "pirenu" is "three dogs"?} \\
\texttt{                  {-} But "tipi" is "three dogs" in S2.} \\
\texttt{                  {-} Maybe "tipi" is "dog" and the number "three" is implied?} \\
\texttt{                  {-} But "a dog" is "tipinu" in S2.} \\
\texttt{                  {-} "tipi" (3 dogs) vs "tipinu" (1 dog).} \\
\texttt{                  {-} This suggests "tipi" is plural and "tipinu" is singular?} \\
\texttt{                  {-} But "tipi" is "three dogs" and "tipinu" is "a dog".} \\
\texttt{                  {-} So "tipi" = plural dog, "tipinu" = singular dog.} \\
\texttt{                  {-} But} \\
            }
        }
    }
    & & \\ \\

    \theutterance \stepcounter{utterance}  
    & & & \multicolumn{2}{p{0.3\linewidth}}{
        \cellcolor[rgb]{0.9,0.9,0.9}{
            \makecell[{{p{\linewidth}}}]{
                \texttt{\tiny{[GM$|$GM]}}
                \texttt{could not find any probe sentences; put them under a 'PROBES:' line, one per line} \\
            }
        }
    }
    & & \\ \\

    \theutterance \stepcounter{utterance}  
    & & & \multicolumn{4}{p{0.6\linewidth}}{
        \cellcolor[rgb]{0.9,0.9,0.9}{
            \makecell[{{p{\linewidth}}}]{
                \texttt{\tiny{[P1$\langle$GM]}}
                \texttt{Your last response could not be processed: could not find any probe sentences; put them under a 'PROBES:' line, one per line.} \\
\texttt{List your probes (a brief THOUGHTS line is optional):} \\
\texttt{PROBES:} \\
\texttt{1. <an English sentence>} \\
\texttt{(up to 3 probe lines; at least one)} \\
            }
        }
    }
     \\ \\

    \theutterance \stepcounter{utterance}  
    & \multicolumn{4}{p{0.6\linewidth}}{
        \cellcolor[rgb]{0.561,0.737,0.561}{
            \makecell[{{p{\linewidth}}}]{
                \texttt{\tiny{[P1$\rangle$GM]}}
                \texttt{THOUGHTS: The previous data shows that singular nouns have no case/number markers (e.g., "dugu" for dog, "tipi" for boat, "ruditi" for tree, "guti" for hill). Plural nouns seem to take a suffix like "{-}nu" or "{-}inu" or "{-}enu", but the mapping is inconsistent (e.g., "three trees" is "dugunu" in one sentence and "dobudinu" in another). This suggests the noun stem itself changes based on number or case, or there is a complex agreement system. I need to isolate the plural forms for specific nouns to see if the stem changes or if the suffix is consistent. I will probe with "two" and "three" for the same noun in different positions (Subject/Object) to see if the form changes. I also need to confirm the verb forms for "chase", "carry", "hear" with singulars.} \\
\\ 
\texttt{PROBES:} \\
\texttt{1. two dogs watch a dog} \\
\texttt{2. three dogs watch a dog} \\
\texttt{3. a dog chase a boat} \\
            }
        }
    }
    & & \\ \\

    \theutterance \stepcounter{utterance}  
    & & & \multicolumn{2}{p{0.3\linewidth}}{
        \cellcolor[rgb]{0.9,0.9,0.9}{
            \makecell[{{p{\linewidth}}}]{
                \texttt{\tiny{[GM$|$GM]}}
                {'kind': 'investigation', 'round': 1, 'probes_issued': 3, 'malformed': 0, 'probes': ['two dogs watch a dog', 'three dogs watch a dog', 'a dog chase a boat']}
            }
        }
    }
    & & \\ \\

    \theutterance \stepcounter{utterance}  
    & & & \multicolumn{4}{p{0.6\linewidth}}{
        \cellcolor[rgb]{0.9,0.9,0.9}{
            \makecell[{{p{\linewidth}}}]{
                \texttt{\tiny{[P1$\langle$GM]}}
                \texttt{TRANSLATIONS:} \\
\texttt{1. "two dogs watch a dog"  =>  tipi deniru tipinu nu} \\
\texttt{2. "three dogs watch a dog"  =>  tipi deniru tipinu ga} \\
\texttt{3. "a dog chase a boat"  =>  dugu buro tipi} \\
\texttt{STATUS: investigation round 2 of 4 done. Probes used: 6 of 12. 2 investigation round(s) left before the exam (10 sentences).} \\
\texttt{Send your next probes now (a PROBES list; brief THOUGHTS optional).} \\
            }
        }
    }
     \\ \\

    \theutterance \stepcounter{utterance}  
    & \multicolumn{4}{p{0.6\linewidth}}{
        \cellcolor[rgb]{0.561,0.737,0.561}{
            \makecell[{{p{\linewidth}}}]{
                \texttt{\tiny{[P1$\rangle$GM]}}
                \texttt{THOUGHTS: The new data is a breakthrough.} \\
\texttt{1. "two dogs watch a dog" => "tipi deniru tipinu nu"} \\
\texttt{2. "three dogs watch a dog" => "tipi deniru tipinu ga"} \\
\texttt{   {-} "tipi" = "two dogs" AND "three dogs"? No, the subject is "tipi" in both.} \\
\texttt{   {-} Wait, "two dogs" and "three dogs" both translate to "tipi"?} \\
\texttt{   {-} Let's re{-}read carefully.} \\
\texttt{   {-} "two dogs watch a dog" => "tipi deniru tipinu nu"} \\
\texttt{   {-} "three dogs watch a dog" => "tipi deniru tipinu ga"} \\
\texttt{   {-} The subject "tipi" is the same for "two dogs" and "three dogs".} \\
\texttt{   {-} The object "tipinu" is the same for "a dog" in both.} \\
\texttt{   {-} The difference is the final particle: "nu" vs "ga".} \\
\texttt{   {-} "nu" corresponds to "two" (in the subject? or object?).} \\
\texttt{   {-} "ga" corresponds to "three" (in the subject? or object?).} \\
\texttt{   {-} But the subject is "tipi" in both. The object is "tipinu" in both.} \\
\texttt{   {-} The only difference is the final word: "nu" vs "ga".} \\
\texttt{   {-} In sentence 1: "two dogs" (S) ... "a dog" (O) ... "nu".} \\
\texttt{   {-} In sentence 2: "three dogs" (S) ... "a dog" (O) ... "ga".} \\
\texttt{   {-} This suggests the final particle marks the number of the **subject**?} \\
\texttt{   {-} Or maybe the number of the **object**?} \\
\texttt{   {-} Let's check the object: "a dog" is singular in both.} \\
\texttt{   {-} So the particle "nu" vs "ga" must mark the number of the **subject**.} \\
\texttt{   {-} "two dogs" {-}> "nu".} \\
\texttt{   {-} "three dogs" {-}> "ga".} \\
\texttt{   {-} But wait, "tipi" is the subject. Does "tipi" change? No, it's "tipi" in both.} \\
\texttt{   {-} So the number is marked by a separate particle at the end of the sentence?} \\
\texttt{   {-} Let's check the previous data with this hypothesis.} \\
\texttt{   {-} "three trees watch two boats" => "dugunu nu deniru gutinu ga".} \\
\texttt{     {-} Subject: "three trees". Object: "two boats".} \\
\texttt{     {-} If the final particle marks the subject number: "ga" should be "three".} \\
\texttt{     {-} But the sentence ends with "ga".} \\
\texttt{     {-} "three trees" {-}> "ga".} \\
\texttt{     {-} "two boats" {-}> "nu"?} \\
\texttt{     {-} But the sentence has "nu" after "dugunu" and "ga" after "gutinu".} \\
\texttt{     {-} "dugunu nu" ... "gutinu ga".} \\
\texttt{     {-} Maybe "nu" marks "two" and "ga" marks "three"?} \\
\texttt{     {-} "dugunu" = "three trees"? "nu" = marker for "three"?} \\
\texttt{     {-} "gutinu" = "two boats"? "ga" = marker for "two"?} \\
\texttt{     {-} But in the new data:} \\
\texttt{       {-} "two dogs" (S) ... "nu" (end).} \\
\texttt{       {-} "three dogs" (S) ... "ga" (end).} \\
\texttt{       {-} This contradicts the previous hypothesis if "nu" = "two" and "ga" = "three".} \\
\texttt{       {-} In new data: "two dogs" {-}> "nu". "three dogs" {-}> "ga".} \\
\texttt{       {-} In old data: "three trees" (S) {-}> "dugunu nu". "two boats" (O) {-}> "gutinu ga".} \\
\texttt{       {-} If "nu" = "two" and "ga" = "three", then:} \\
\texttt{         {-} "three trees" (S) should have "ga". But it has "nu".} \\
\texttt{         {-} "two boats" (O) should have "nu". But it has "ga".} \\
\texttt{       {-} This suggests the markers are attached to the noun, not at the end.} \\
\texttt{       {-} "dugunu" = "three trees". "nu" is part of the word?} \\
\texttt{       {-} "gutinu" = "two boats". "ga" is part of the word?} \\
\texttt{       {-} But in new data: "tipi" (S) ... "tipinu" (O) ... "nu" (end).} \\
\texttt{       {-} "tipi" (S) ... "tipinu" (O) ... "ga" (end).} \\
\texttt{       {-} The subject "tipi" is the same. The object "tipinu" is the same.} \\
\texttt{       {-} The only difference is the final particle.} \\
\texttt{       {-} This implies the number of the subject is marked by a final particle?} \\
\texttt{       {-} But "tipi" is "two dogs" and "three dogs".} \\
\texttt{       {-} This means "tipi" is just "dogs" (plural)? And the number is marked by the final particle?} \\
\texttt{       {-} "two dogs" {-}> "tipi ... nu".} \\
\texttt{       {-} "three dogs" {-}> "tipi ... ga".} \\
\texttt{       {-} So "nu" = "two", "ga" = "three".} \\
\texttt{       {-} Let's re{-}check the old data with this.} \\
\texttt{       {-} "three trees watch two boats" => "dugunu nu deniru gutinu ga".} \\
\texttt{         {-} Subject: "three trees". Should end with "ga" (if "ga" = three).} \\
\texttt{         {-} But it ends with "ga" after "gutinu".} \\
\texttt{         {-} "dugunu nu" ... "gutinu ga".} \\
\texttt{         {-} Maybe the structure is {[}Subject{]} {[}Marker{]} {[}Verb{]} {[}Object{]} {[}Marker{]}?} \\
\texttt{         {-} "dugunu" = "trees"? "nu" = "three"?} \\
\texttt{         {-} "gutinu" = "boats"? "ga" = "two"?} \\
\texttt{         {-} But "nu" = "two" and "ga" = "three" from new data.} \\
\texttt{         {-} Contradiction.} \\
\texttt{       {-} Let's re{-}examine the new data.} \\
\texttt{       {-} "two dogs watch a dog" => "tipi deniru tipinu nu".} \\
\texttt{       {-} "three dogs watch a dog" => "tipi deniru tipinu ga".} \\
\texttt{       {-} Subject: "tipi". Object: "tipinu".} \\
\texttt{       {-} "tipi" = "dogs" (plural). "tipinu" = "dog" (singular).} \\
\texttt{       {-} The final particle "nu" vs "ga" marks the number of the **subject**.} \\
\texttt{       {-} "nu" = "two". "ga" = "three".} \\
\texttt{       {-} Now let's look at "three trees watch two boats" => "dugunu nu deniru gutinu ga".} \\
\texttt{         {-} Subject: "three trees". Should have "ga" at the end?} \\
\texttt{         {-} But the sentence ends with "ga".} \\
\texttt{         {-} "dugunu nu" ... "gutinu ga".} \\
\texttt{         {-} Maybe "dugunu" = "trees" (plural)? "nu" = "three"?} \\
\texttt{         {-} "gutinu" = "boats" (plural)? "ga" = "two"?} \\
\texttt{         {-} But "nu" = "two" and "ga" = "three" from new data.} \\
\texttt{         {-} This is a contradiction.} \\
\texttt{       {-} Unless the markers are attached to the noun, and the final particle is something else.} \\
\texttt{       {-} Let's look at "two dogs chase two boats" => "dugunu nu buro tipinu nu".} \\
\texttt{         {-} Subject: "two dogs". Object: "two boats".} \\
\texttt{         {-} "dugunu" = "dogs"? "nu" = "two"?} \\
\texttt{         {-} "tipinu" = "boats"? "nu" = "two"?} \\
\texttt{         {-} If "nu" = "two", then "dugunu" = "dogs" (plural) and "tipinu" = "boats" (plural).} \\
\texttt{         {-} But "tipinu" was "a dog" (singular) in new data.} \\
\texttt{         {-} "a dog" = "tipinu".} \\
\texttt{         {-} "two boats" = "tipinu"?} \\
\texttt{         {-} This implies "tipinu" is "dog" (singular) and "boats" (plural)?} \\
\texttt{         {-} No, "tipi" = "dog" (singular) in new data?} \\
\texttt{         {-} "a dog chase a boat" => "dugu buro tipi".} \\
\texttt{           {-} "dugu" = dog. "tipi" = boat.} \\
\texttt{         {-} "two dogs watch a dog" => "tipi deniru tipinu nu".} \\
\texttt{           {-} "tipi" = "two dogs"? "tipinu" = "a dog"?} \\
\texttt{           {-} But "tipi" was "boat" in "a dog chase a boat".} \\
\texttt{           {-} Contradiction: "tipi" cannot be "boat" and "two dogs".} \\
\texttt{           {-} Unless "tipi" is "dog" and "dugu" is "boat"?} \\
\texttt{           {-} "a dog chase a boat" => "dugu buro tipi".} \\
\texttt{             {-} If "dugu" = boat, "tipi" = dog.} \\
\texttt{             {-} "two dogs watch a dog" => "tipi deniru tipinu nu".} \\
\texttt{               {-} "tipi" = "two dogs". "tipinu" = "a dog".} \\
\texttt{               {-} This works if "tipi" = "dog" (plural) and "tipinu" = "dog" (singular).} \\
\texttt{               {-} But "tipi" is "dog" (plural) and "tipi" is "dog" (singular)?} \\
\texttt{               {-} No, "tipi" is "two dogs" (plural) and "tipinu" is "a dog" (singular).} \\
\texttt{               {-} So "tipi" = "dog" (plural). "tipinu" = "dog" (singular).} \\
\texttt{               {-} But "a dog chase a boat" => "dugu buro tipi".} \\
\texttt{                 {-} "dugu" = dog? "tipi" = boat?} \\
\texttt{                 {-} If "tipi" = "dog" (plural), then "a dog chase a boat" => "dugu buro tipi" means "dog chase dog"?} \\
\texttt{                 {-} No, the English is "a dog chase a boat".} \\
\texttt{                 {-} So "dugu" = dog, "tipi" = boat.} \\
\texttt{                 {-} But "two dogs watch a dog" => "tipi deniru tipinu nu".} \\
\texttt{                   {-} "tipi" = "two dogs".} \\
\texttt{                   {-} "tipinu" = "a dog".} \\
\texttt{                   {-} This implies "tipi" = "dog" (plural) and "tipinu" = "dog" (singular).} \\
\texttt{                   {-} But "tipi" was "boat" in the other sentence.} \\
\texttt{                   {-} This is a contradiction.} \\
\texttt{       {-} Let's re{-}read the new data carefully.} \\
\texttt{       {-} "a dog chase a boat" => "dugu buro tipi".} \\
\texttt{         {-} "dugu" = dog. "buro" = chase. "tipi" = boat.} \\
\texttt{       {-} "two dogs watch a dog" => "tipi deniru tipinu nu".} \\
\texttt{         {-} "tipi" = "two dogs". "deniru" = watch. "tipinu" = "a dog". "nu" = marker.} \\
\texttt{         {-} But "tipi" was "boat" in the previous sentence.} \\
\texttt{         {-} This implies "tipi" is both "boat" and "two dogs".} \\
\texttt{         {-} This is impossible.} \\
\texttt{       {-} Unless "tipi" is "dog" and "dugu" is "boat"?} \\
\texttt{         {-} "a dog chase a boat" => "dugu buro tipi".} \\
\texttt{           {-} If "dugu" = boat, "tipi" = dog.} \\
\texttt{           {-} "two dogs watch a dog" => "tipi deniru tipinu nu".} \\
\texttt{             {-} "tipi" = "two dogs". "tipinu" = "a dog".} \\
\texttt{             {-} This works if "tipi" = "dog" (plural) and "tipinu" = "dog" (singular).} \\
\texttt{             {-} But "tipi" is "dog" (plural) and "tipi" is "dog" (singular)?} \\
\texttt{             {-} No, "tipi" is "two dogs" (plural) and "tipinu" is "a dog" (singular).} \\
\texttt{             {-} So "tipi" = "dog" (plural). "tipinu" = "dog" (singular).} \\
\texttt{             {-} But "a dog chase a boat" => "dugu buro tipi".} \\
\texttt{               {-} "dugu" = boat. "tipi" = dog.} \\
\texttt{               {-} This works.} \\
\texttt{             {-} So "dugu" = boat. "tipi" = dog.} \\
\texttt{             {-} But earlier I thought "dugu" = dog.} \\
\texttt{             {-} Let's check "a dog watch a boat" => "dugu deniru tipi".} \\
\texttt{               {-} "dugu" = dog. "tipi" = boat.} \\
\texttt{               {-} This contradicts "dugu" = boat.} \\
\texttt{             {-} So "dugu" = dog, "tipi" = boat.} \\
\texttt{             {-} Then "two dogs watch a dog" => "tipi deniru tipinu nu".} \\
\texttt{               {-} "tipi" = "two dogs".} \\
\texttt{               {-} But "tipi" = "boat".} \\
\texttt{               {-} Contradiction.} \\
\texttt{       {-} Maybe the translation of "two dogs watch a dog" is "tipi deniru tipinu nu" where "tipi" = "two dogs" and "tipinu" = "a dog".} \\
\texttt{       {-} And "a dog watch a boat" => "dugu deniru tipi" where "dugu" = "a dog" and "tipi" = "a boat".} \\
\texttt{       {-} This implies "tipi" is "boat" and "two dogs".} \\
\texttt{       {-} This is impossible.} \\
\texttt{       {-} Unless "tipi" is "dog" and "dugu" is "boat"?} \\
\texttt{         {-} "a dog watch a boat" => "dugu deniru tipi".} \\
\texttt{           {-} "dugu" = dog. "tipi" = boat.} \\
\texttt{         {-} "two dogs watch a dog" => "tipi deniru tipinu nu".} \\
\texttt{           {-} "tipi" = "two dogs". "tipinu" = "a dog".} \\
\texttt{           {-} This implies "tipi" = "dog" (plural) and "tipinu" = "dog" (singular).} \\
\texttt{           {-} But "tipi" was "boat" in the first sentence.} \\
\texttt{           {-} Contradiction.} \\
\texttt{       {-} Maybe the informant made a mistake? No, the informant is truthful.} \\
\texttt{       {-} Maybe I am mis{-}segmenting.} \\
\texttt{       {-} "two dogs watch a dog" => "tipi deniru tipinu nu".} \\
\texttt{         {-} Maybe "tipi" = "two", "deniru" = "dogs watch", "tipinu" = "a dog", "nu" = marker?} \\
\texttt{         {-} No, "deniru" is "watch" from "a dog watch a boat" => "dugu deniru tipi".} \\
\texttt{         {-} So "deniru" = "watch".} \\
\texttt{         {-} "tipi" = "two dogs". "tipinu" = "a dog".} \\
\texttt{         {-} "tipi" = "two dogs". "tipinu" = "a dog".} \\
\texttt{         {-} "tipi" and "tipinu" are related.} \\
\texttt{         {-} "tipi" = "dog" (plural). "tipinu" = "dog" (singular).} \\
\texttt{         {-} But "a dog watch a boat" => "dugu deniru tipi".} \\
\texttt{           {-} "dugu" = "a dog". "tipi" = "a boat".} \\
\texttt{           {-} So "tipi" = "boat".} \\
\texttt{           {-} Contradiction: "tipi" cannot be "dog" (plural) and "boat".} \\
\texttt{       {-} Unless "tipi" is "dog" and "dugu" is "boat"?} \\
\texttt{         {-} "a dog watch a boat" => "dugu deniru tipi".} \\
\texttt{           {-} "dugu" = dog. "tipi" = boat.} \\
\texttt{         {-} "two dogs watch a dog" => "tipi deniru tipinu nu".} \\
\texttt{           {-} "tipi" = "two dogs". "tipinu" = "a dog".} \\
\texttt{           {-} This implies "tipi" = "dog" (plural).} \\
\texttt{           {-} But "tipi" = "boat".} \\
\texttt{           {-} Contradiction.} \\
\texttt{       {-} Maybe the word for "dog" is "tipi" and "boat" is "dugu"?} \\
\texttt{         {-} "a dog watch a boat" => "dugu deniru tipi".} \\
\texttt{           {-} "dugu" = dog. "tipi" = boat.} \\
\texttt{         {-} "two dogs watch a dog" => "tipi deniru tipinu nu".} \\
\texttt{           {-} "tipi" = "two dogs". "tipinu" = "a dog".} \\
\texttt{           {-} This implies "tipi" = "dog" (plural).} \\
\texttt{           {-} But "tipi" = "boat".} \\
\texttt{           {-} Contradiction.} \\
\texttt{       {-} Maybe the translation of "two dogs watch a dog" is "tipi deniru tipinu nu" where "tipi" = "two", "deniru" = "dogs watch", "tipinu" = "a dog", "nu" = marker?} \\
\texttt{         {-} No, "deniru" is "watch".} \\
\texttt{         {-} "tipi" = "two dogs". "tipinu" = "a dog".} \\
\texttt{         {-} "tipi" = "two dogs". "tipinu" = "a dog".} \\
\texttt{         {-} "tipi" and "tipinu" are related.} \\
\texttt{         {-} "tipi" = "dog" (plural). "tipinu" = "dog" (singular).} \\
\texttt{         {-} But "a dog watch a boat" => "dugu deniru tipi".} \\
\texttt{           {-} "dugu" = "a dog". "tipi" = "a boat".} \\
\texttt{           {-} So "tipi" = "boat".} \\
\texttt{           {-} Contradiction.} \\
\texttt{       {-} Maybe the word for "dog" is "dugu" and "boat" is "tipi"?} \\
\texttt{         {-} "a dog watch a boat" => "dugu deniru tipi".} \\
\texttt{           {-} "dugu" = dog. "tipi" = boat.} \\
\texttt{         {-} "two dogs watch a dog" => "tipi deniru tipinu nu".} \\
\texttt{           {-} "tipi" = "two dogs". "tipinu" = "a dog".} \\
\texttt{           {-} This implies "tipi" = "dog" (plural).} \\
\texttt{           {-} But "tipi" = "boat".} \\
\texttt{           {-} Contradiction.} \\
\texttt{       {-} Maybe the word for "dog" is "tipi" and "boat" is "dugu"?} \\
\texttt{         {-} "a dog watch a boat" => "dugu deniru tipi".} \\
\texttt{           {-} "dugu" = dog. "tipi" = boat.} \\
\texttt{         {-} "two dogs watch a dog" => "tipi deniru tipinu nu".} \\
\texttt{           {-} "tipi" = "two dogs". "tipinu" = "a dog".} \\
\texttt{           {-} This implies "tipi" = "dog" (plural).} \\
\texttt{           {-} But "tipi" = "boat".} \\
\texttt{           {-} Contradiction.} \\
\texttt{       {-} Maybe the translation of "two dogs watch a dog" is "tipi deniru tipinu nu" where "tipi" = "two", "deniru" = "dogs watch", "tipinu" = "a dog", "nu" = marker?} \\
\texttt{         {-} No, "deniru" is "watch".} \\
\texttt{         {-} "tipi" = "two dogs". "tipinu" = "a dog".} \\
\texttt{         {-} "tipi" = "two dogs". "tipinu" = "a dog".} \\
\texttt{         {-} "tipi" and "tipinu" are related.} \\
\texttt{         {-} "tipi" = "dog" (plural). "tipinu" = "dog" (singular).} \\
\texttt{         {-} But "a dog watch a boat" => "dugu deniru tipi".} \\
\texttt{           {-} "dugu" = "a dog". "tipi" = "a boat".} \\
\texttt{           {-} So "tipi" = "boat".} \\
\texttt{           {-} Contradiction.} \\
\texttt{       {-} Maybe the word for "dog" is "dugu" and "boat" is "tipi"?} \\
\texttt{         {-} "a dog watch a boat" => "dugu deniru tipi".} \\
\texttt{           {-} "dugu" = dog. "tipi" = boat.} \\
\texttt{         {-} "two dogs watch a dog" => "tipi deniru tipinu nu".} \\
\texttt{           {-} "tipi" = "two dogs". "tipinu" = "a dog".} \\
\texttt{           {-} This implies "tipi" = "dog" (plural).} \\
\texttt{           {-} But "tipi" = "boat".} \\
\texttt{           {-} Contradiction.} \\
\texttt{       {-} Maybe the word for "dog" is "tipi" and "boat" is "dugu"?} \\
\texttt{         {-} "a dog watch a boat" => "dugu deniru tipi".} \\
\texttt{           {-} "dugu" = dog. "tipi" = boat.} \\
\texttt{         {-} "two dogs watch a dog" => "tipi deniru tipinu nu".} \\
\texttt{           {-} "tipi" = "two dogs". "tipinu" = "a dog".} \\
\texttt{           {-} This implies "tipi" = "dog" (plural).} \\
\texttt{           {-} But "tipi" = "boat".} \\
\texttt{           {-} Contradiction.} \\
\texttt{       {-} Maybe the translation of "two dogs watch a dog" is "tipi deniru tipinu nu" where "tipi" = "two", "deniru" = "dogs watch", "tipinu" = "a dog", "nu" = marker?} \\
\texttt{         {-} No, "deniru" is "watch".} \\
\texttt{         {-} "tipi" = "two dogs". "tipinu" = "a dog".} \\
\texttt{         {-} "tipi" = "two dogs". "tipinu" = "a dog".} \\
\texttt{         {-} "tipi" and "tipinu" are related.} \\
\texttt{         {-} "tipi" = "dog" (plural). "tipinu" = "dog" (singular).} \\
\texttt{         {-} But "a dog watch a boat" => "dugu deniru tipi".} \\
\texttt{           {-} "dugu" = "a dog". "tipi" = "a boat".} \\
\texttt{           {-} So "tipi" = "boat".} \\
\texttt{           {-} Contradiction.} \\
\texttt{       {-} Maybe the word for "dog" is "dugu" and "boat" is "tipi"?} \\
\texttt{         {-} "a dog watch a boat" => "dugu deniru tipi".} \\
\texttt{           {-} "dugu" = dog. "tipi" = boat.} \\
\texttt{         {-} "two dogs watch a dog" => "tipi deniru tipinu nu".} \\
\texttt{           {-} "tipi" = "two dogs". "tipinu" = "a dog".} \\
\texttt{           {-} This implies "tipi" = "dog" (plural).} \\
\texttt{           {-} But "tipi"} \\
            }
        }
    }
    & & \\ \\

    \theutterance \stepcounter{utterance}  
    & & & \multicolumn{2}{p{0.3\linewidth}}{
        \cellcolor[rgb]{0.9,0.9,0.9}{
            \makecell[{{p{\linewidth}}}]{
                \texttt{\tiny{[GM$|$GM]}}
                \texttt{could not find any probe sentences; put them under a 'PROBES:' line, one per line} \\
            }
        }
    }
    & & \\ \\

    \theutterance \stepcounter{utterance}  
    & & & \multicolumn{4}{p{0.6\linewidth}}{
        \cellcolor[rgb]{0.9,0.9,0.9}{
            \makecell[{{p{\linewidth}}}]{
                \texttt{\tiny{[P1$\langle$GM]}}
                \texttt{Your last response could not be processed: could not find any probe sentences; put them under a 'PROBES:' line, one per line.} \\
\texttt{List your probes (a brief THOUGHTS line is optional):} \\
\texttt{PROBES:} \\
\texttt{1. <an English sentence>} \\
\texttt{(up to 3 probe lines; at least one)} \\
            }
        }
    }
     \\ \\

    \theutterance \stepcounter{utterance}  
    & \multicolumn{4}{p{0.6\linewidth}}{
        \cellcolor[rgb]{0.561,0.737,0.561}{
            \makecell[{{p{\linewidth}}}]{
                \texttt{\tiny{[P1$\rangle$GM]}}
                \texttt{THOUGHTS: The previous data revealed a critical contradiction: "tipi" appears to mean "boat" in "a dog chase a boat" (dugu buro tipi) but "two dogs" in "two dogs watch a dog" (tipi deniru tipinu nu). This suggests my segmentation or understanding of the noun stems is wrong.} \\
\texttt{Hypothesis: The noun stem changes based on number or case, or "tipi" is a number word?} \\
\texttt{Let's re{-}examine "a dog chase a boat" => "dugu buro tipi".} \\
\texttt{If "dugu" = dog, "buro" = chase, "tipi" = boat.} \\
\texttt{"two dogs watch a dog" => "tipi deniru tipinu nu".} \\
\texttt{If "tipi" = two dogs, then "tipi" cannot be boat.} \\
\texttt{Unless "tipi" is "dog" and "dugu" is "boat"?} \\
\texttt{"a dog chase a boat" => "dugu buro tipi". If "dugu" = boat, "tipi" = dog.} \\
\texttt{"two dogs watch a dog" => "tipi deniru tipinu nu". "tipi" = two dogs. "tipinu" = a dog.} \\
\texttt{This works if "tipi" = dog (plural) and "tipinu" = dog (singular).} \\
\texttt{But "a dog chase a boat" => "dugu buro tipi". If "tipi" = dog, then "dugu" = boat.} \\
\texttt{Let's check "a dog watch a boat" => "dugu deniru tipi".} \\
\texttt{If "dugu" = boat, "tipi" = dog.} \\
\texttt{Then "a dog watch a boat" means "boat watch dog"? No, English is "a dog watch a boat".} \\
\texttt{So "dugu" must be dog, "tipi" must be boat.} \\
\texttt{Then "two dogs watch a dog" => "tipi deniru tipinu nu" means "boat watch dog"? No.} \\
\texttt{This implies "tipi" in "two dogs watch a dog" is NOT "boat".} \\
\texttt{Maybe "tipi" is "two"?} \\
\texttt{"two dogs" = "tipi" + "dog"?} \\
\texttt{But "tipi" is the first word.} \\
\texttt{Maybe the word order is different?} \\
\texttt{Or maybe "tipi" is "dog" and "dugu" is "boat"?} \\
\texttt{Let's try: "dugu" = boat, "tipi" = dog.} \\
\texttt{"a dog watch a boat" => "dugu deniru tipi".} \\
\texttt{If "dugu" = boat, "tipi" = dog. Then "boat watch dog". Incorrect.} \\
\texttt{So "dugu" = dog, "tipi" = boat.} \\
\texttt{Then "two dogs watch a dog" => "tipi deniru tipinu nu".} \\
\texttt{"tipi" = two dogs? But "tipi" = boat.} \\
\texttt{Contradiction.} \\
\texttt{Unless "tipi" is a number word "two"?} \\
\texttt{"two dogs" = "tipi" (two) + "dog" (stem)?} \\
\texttt{But where is the stem for dog?} \\
\texttt{Maybe "tipi" is "dog" and "dugu" is "boat"?} \\
\texttt{Let's assume the informant made a mistake? No.} \\
\texttt{Maybe "tipi" is "dog" and "dugu" is "boat" and the translation "a dog watch a boat" is "dugu deniru tipi" meaning "boat watch dog"? No, the English is fixed.} \\
\texttt{Wait, maybe "dugu" is "dog" and "tipi" is "boat" is correct.} \\
\texttt{Then "two dogs watch a dog" => "tipi deniru tipinu nu".} \\
\texttt{"tipi" = two dogs.} \\
\texttt{This implies "tipi" is a word for "two dogs".} \\
\texttt{But "tipi" is also "boat".} \\
\texttt{This is impossible unless "boat" and "dog" are the same word? No.} \\
\texttt{Maybe "tipi" is "two" and the noun is omitted? No.} \\
\texttt{Let's try a different angle.} \\
\texttt{Maybe "tipi" is "dog" and "dugu" is "boat" and the translation "a dog watch a boat" is actually "dugu deniru tipi" where "dugu" = dog and "tipi" = boat.} \\
\texttt{Then "two dogs watch a dog" => "tipi deniru tipinu nu".} \\
\texttt{"tipi" = two dogs.} \\
\texttt{This implies "tipi" = "two dogs".} \\
\texttt{But "tipi" = "boat".} \\
\texttt{Contradiction.} \\
\texttt{Unless "tipi" is "dog" and "dugu" is "boat" and the translation "a dog watch a boat" is "dugu deniru tipi" where "dugu" = boat and "tipi" = dog?} \\
\texttt{Then "boat watch dog". Incorrect.} \\
\texttt{Maybe the word order is Object{-}Verb{-}Subject?} \\
\texttt{"a dog watch a boat" => "dugu deniru tipi".} \\
\texttt{If O{-}V{-}S: "dugu" (boat) "deniru" (watch) "tipi" (dog).} \\
\texttt{"boat watch dog". Incorrect.} \\
\texttt{If S{-}V{-}O: "dugu" (dog) "deniru" (watch) "tipi" (boat). Correct.} \\
\texttt{So S{-}V{-}O is correct.} \\
\texttt{Then "two dogs watch a dog" => "tipi deniru tipinu nu".} \\
\texttt{S{-}V{-}O: "tipi" (two dogs) "deniru" (watch) "tipinu" (a dog).} \\
\texttt{So "tipi" = two dogs.} \\
\texttt{But "tipi" = boat.} \\
\texttt{Contradiction.} \\
\texttt{Maybe "tipi" is "dog" and "dugu" is "boat" and the translation "a dog watch a boat" is "dugu deniru tipi" where "dugu" = dog and "tipi" = boat.} \\
\texttt{Then "two dogs watch a dog" => "tipi deniru tipinu nu".} \\
\texttt{"tipi" = two dogs.} \\
\texttt{This implies "tipi" = "two dogs".} \\
\texttt{But "tipi" = "boat".} \\
\texttt{This is a paradox.} \\
\texttt{Unless "tipi" is "dog" and "dugu" is "boat" and the translation "a dog watch a boat" is "dugu deniru tipi" where "dugu" = boat and "tipi" = dog?} \\
\texttt{Then "boat watch dog". Incorrect.} \\
\texttt{Maybe the word for "dog" is "tipi" and "boat" is "dugu"?} \\
\texttt{"a dog watch a boat" => "dugu deniru tipi".} \\
\texttt{If "dugu" = dog, "tipi" = boat.} \\
\texttt{"two dogs watch a dog" => "tipi deniru tipinu nu".} \\
\texttt{"tipi" = two dogs.} \\
\texttt{This implies "tipi" = "two dogs".} \\
\texttt{But "tipi" = "boat".} \\
\texttt{Contradiction.} \\
\texttt{Maybe "tipi" is "two" and "dog" is "tipi"?} \\
\texttt{"two dogs" = "tipi" (two) + "tipi" (dog)? No.} \\
\texttt{Maybe "tipi" is "dog" and "dugu" is "boat" and the translation "a dog watch a boat" is "dugu deniru tipi" where "dugu" = dog and "tipi" = boat.} \\
\texttt{Then "two dogs watch a dog" => "tipi deniru tipinu nu".} \\
\texttt{"tipi" = two dogs.} \\
\texttt{This implies "tipi" = "two dogs".} \\
\texttt{But "tipi" = "boat".} \\
\texttt{This is impossible.} \\
\texttt{Unless "tipi" is "dog" and "dugu" is "boat" and the translation "a dog watch a boat" is "dugu deniru tipi" where "dugu" = boat and "tipi" = dog?} \\
\texttt{Then "boat watch dog". Incorrect.} \\
\texttt{Maybe the word for "dog" is "dugu" and "boat" is "tipi" is correct.} \\
\texttt{Then "two dogs watch a dog" => "tipi deniru tipinu nu".} \\
\texttt{"tipi" = two dogs.} \\
\texttt{This implies "tipi" = "two dogs".} \\
\texttt{But "tipi" = "boat".} \\
\texttt{This suggests "tipi" is a homonym? Unlikely.} \\
\texttt{Maybe "tipi" is "dog" and "dugu" is "boat" and the translation "a dog watch a boat" is "dugu deniru tipi" where "dugu" = dog and "tipi" = boat.} \\
\texttt{Then "two dogs watch a dog" => "tipi deniru tipinu nu".} \\
\texttt{"tipi" = two dogs.} \\
\texttt{This implies "tipi" = "two dogs".} \\
\texttt{But "tipi" = "boat".} \\
\texttt{This is a contradiction.} \\
\texttt{Maybe the translation of "two dogs watch a dog" is "tipi deniru tipinu nu" where "tipi" = "two", "deniru" = "dogs watch", "tipinu" = "a dog", "nu" = marker?} \\
\texttt{No, "deniru" is "watch".} \\
\texttt{Maybe "tipi" is "dog" and "dugu" is "boat" and the translation "a dog watch a boat" is "dugu deniru tipi" where "dugu" = dog and "tipi" = boat.} \\
\texttt{Then "two dogs watch a dog" => "tipi deniru tipinu nu".} \\
\texttt{"tipi" = two dogs.} \\
\texttt{This implies "tipi" = "two dogs".} \\
\texttt{But "tipi" = "boat".} \\
\texttt{This is impossible.} \\
\texttt{Maybe the word for "dog" is "tipi" and "boat" is "dugu"?} \\
\texttt{"a dog watch a boat" => "dugu deniru tipi".} \\
\texttt{If "dugu" = dog, "tipi" = boat.} \\
\texttt{"two dogs watch a dog" => "tipi deniru tipinu nu".} \\
\texttt{"tipi" = two dogs.} \\
\texttt{This implies "tipi" = "two dogs".} \\
\texttt{But "tipi" = "boat".} \\
\texttt{Contradiction.} \\
\texttt{Maybe "tipi" is "dog" and "dugu" is "boat" and the translation "a dog watch a boat" is "dugu deniru tipi" where "dugu" = boat and "tipi" = dog?} \\
\texttt{Then "boat watch dog". Incorrect.} \\
\texttt{Maybe the word for "dog" is "dugu" and "boat" is "tipi" is correct.} \\
\texttt{Then "two dogs watch a dog" => "tipi deniru tipinu nu".} \\
\texttt{"tipi" = two dogs.} \\
\texttt{This implies "tipi" = "two dogs".} \\
\texttt{But "tipi" = "boat".} \\
\texttt{This suggests "tipi" is a number word "two"?} \\
\texttt{"two dogs" = "tipi" (two) + "dog" (stem)?} \\
\texttt{But where is the stem for dog?} \\
\texttt{Maybe "tipi" is "dog" and "dugu" is "boat" and the translation "a dog watch a boat" is "dugu deniru tipi" where "dugu" = dog and "tipi" = boat.} \\
\texttt{Then "two dogs watch a dog" => "tipi deniru tipinu nu".} \\
\texttt{"tipi" = two dogs.} \\
\texttt{This implies "tipi" = "two dogs".} \\
\texttt{But "tipi" = "boat".} \\
\texttt{This is a paradox.} \\
\texttt{Maybe the word for "dog" is "tipi" and "boat" is "dugu"?} \\
\texttt{"a dog watch a boat" => "dugu deniru tipi".} \\
\texttt{If "dugu" = dog, "tipi" = boat.} \\
\texttt{"two dogs watch a dog" => "tipi deniru tipinu nu".} \\
\texttt{"tipi" = two dogs.} \\
\texttt{This implies "tipi" = "two dogs".} \\
\texttt{But "tipi" = "boat".} \\
\texttt{Contradiction.} \\
\texttt{Maybe "tipi" is "dog" and "dugu" is "boat" and the translation "a dog watch a boat" is "dugu deniru tipi" where "dugu" = boat and "tipi" = dog?} \\
\texttt{Then "boat watch dog". Incorrect.} \\
\texttt{Maybe the word for "dog" is "dugu" and "boat" is "tipi" is correct.} \\
\texttt{Then "two dogs watch a dog" => "tipi deniru tipinu nu".} \\
\texttt{"tipi" = two dogs.} \\
\texttt{This implies "tipi" = "two dogs".} \\
\texttt{But "tipi" = "boat".} \\
\texttt{This suggests "tipi" is a homonym? Unlikely.} \\
\texttt{Maybe "tipi" is "dog" and "dugu" is "boat" and the translation "a dog watch a boat" is "dugu deniru tipi" where "dugu" = dog and "tipi" = boat.} \\
\texttt{Then "two dogs watch a dog" => "tipi deniru tipinu nu".} \\
\texttt{"tipi" = two dogs.} \\
\texttt{This implies "tipi" = "two dogs".} \\
\texttt{But "tipi" = "boat".} \\
\texttt{This is impossible.} \\
\texttt{Maybe the translation of "two dogs watch a dog" is "tipi deniru tipinu nu" where "tipi" = "two", "deniru" = "dogs watch", "tipinu" = "a dog", "nu" = marker?} \\
\texttt{No, "deniru" is "watch".} \\
\texttt{Maybe "tipi" is "dog" and "dugu" is "boat" and the translation "a dog watch a boat" is "dugu deniru tipi" where "dugu" = dog and "tipi" = boat.} \\
\texttt{Then "two dogs watch a dog" => "tipi deniru tipinu nu".} \\
\texttt{"tipi" = two dogs.} \\
\texttt{This implies "tipi" = "two dogs".} \\
\texttt{But "tipi" = "boat".} \\
\texttt{This is a contradiction.} \\
\texttt{Maybe the word for "dog" is "tipi" and "boat" is "dugu"?} \\
\texttt{"a dog watch a boat" => "dugu deniru tipi".} \\
\texttt{If "dugu" = dog, "tipi" = boat.} \\
\texttt{"two dogs watch a dog" => "tipi deniru tipinu nu".} \\
\texttt{"tipi" = two dogs.} \\
\texttt{This implies "tipi" = "two dogs".} \\
\texttt{But "tipi" = "boat".} \\
\texttt{Contradiction.} \\
\texttt{Maybe "tipi" is "dog" and "dugu" is "boat" and the translation "a dog watch a boat" is "dugu deniru tipi" where "dugu" = boat and "tipi" = dog?} \\
\texttt{Then "boat watch dog". Incorrect.} \\
\texttt{Maybe the word for "dog" is "dugu" and "boat" is "tipi" is correct.} \\
\texttt{Then "two dogs watch a dog" => "tipi deniru tipinu nu".} \\
\texttt{"tipi" = two dogs.} \\
\texttt{This implies "tipi" = "two dogs".} \\
\texttt{But "tipi" = "boat".} \\
\texttt{This suggests "tipi" is a number word "two"?} \\
\texttt{"two dogs" = "tipi" (two) + "dog" (stem)?} \\
\texttt{But where is the stem for dog?} \\
\texttt{Maybe "tipi" is "dog" and "dugu" is "boat" and the translation "a dog watch a boat" is "dugu deniru tipi" where "dugu" = dog and "tipi" = boat.} \\
\texttt{Then "two dogs watch a dog" => "tipi deniru tipinu nu".} \\
\texttt{"tipi" = two dogs.} \\
\texttt{This implies "tipi" = "two dogs".} \\
\texttt{But "tipi" = "boat".} \\
\texttt{This is a paradox.} \\
\texttt{Maybe the word for "dog" is "tipi" and "boat" is "dugu"?} \\
\texttt{"a dog watch a boat" => "dugu deniru tipi".} \\
\texttt{If "dugu" = dog, "tipi" = boat.} \\
\texttt{"two dogs watch a dog" => "tipi deniru tipinu nu".} \\
\texttt{"tipi" = two dogs.} \\
\texttt{This implies "tipi" = "two dogs".} \\
\texttt{But "tipi" = "boat".} \\
\texttt{Contradiction.} \\
\texttt{Maybe "tipi" is "dog" and "dugu" is "boat" and the translation "a dog watch a boat" is "dugu deniru tipi" where "dugu" = boat and "tipi" = dog?} \\
\texttt{Then "boat watch dog". Incorrect.} \\
\texttt{Maybe the word for "dog" is "dugu" and "boat" is "tipi" is correct.} \\
\texttt{Then "two dogs watch a dog" => "tipi deniru tipinu nu".} \\
\texttt{"tipi" = two dogs.} \\
\texttt{This implies "tipi" = "two dogs".} \\
\texttt{But "tipi" = "boat".} \\
\texttt{This suggests "tipi" is a homonym? Unlikely.} \\
\texttt{Maybe "tipi" is "dog" and "dugu" is "boat" and the translation "a dog watch a boat" is "dugu deniru tipi" where "dugu" = dog and "tipi" = boat.} \\
\texttt{Then "two dogs watch a dog" => "tipi deniru tipinu nu".} \\
\texttt{"tipi" = two dogs.} \\
\texttt{This implies "tipi" = "two dogs".} \\
\texttt{But "tipi" = "boat".} \\
\texttt{This is impossible.} \\
\texttt{Maybe the translation of "two dogs watch a dog" is "tipi deniru tipinu nu" where "tipi" = "two", "deniru" = "dogs watch", "tipinu" = "a dog", "nu" = marker?} \\
\texttt{No, "deniru" is "watch".} \\
\texttt{Maybe "tipi" is "dog" and "dugu" is "boat" and the translation "a dog watch a boat" is "dugu deniru tipi" where "dugu" = dog and "tipi" = boat.} \\
\texttt{Then "two dogs watch a dog" => "tipi deniru tipinu nu".} \\
\texttt{"tipi" = two dogs.} \\
\texttt{This implies "tipi" = "two dogs".} \\
\texttt{But "tipi" = "boat".} \\
\texttt{This is a contradiction.} \\
\texttt{Maybe the word for "dog" is "tipi" and "boat" is "dugu"?} \\
\texttt{"a dog watch a boat" => "dugu deniru tipi".} \\
\texttt{If "dugu" = dog, "tipi" = boat.} \\
\texttt{"two dogs watch a dog" => "tipi deniru tipinu nu".} \\
\texttt{"tipi" = two dogs.} \\
\texttt{This implies "tipi" = "two dogs".} \\
\texttt{But "tipi" = "boat".} \\
\texttt{Contradiction.} \\
\texttt{Maybe "tipi" is "dog" and "dugu" is "boat" and the translation "a dog watch a boat" is "dugu deniru tipi" where "dugu" = boat and "tipi" = dog?} \\
\texttt{Then "boat watch dog". Incorrect.} \\
\texttt{Maybe the word for "dog" is "dugu" and "boat" is "tipi" is correct.} \\
\texttt{Then "two dogs watch a dog" => "tipi deniru tipinu nu".} \\
\texttt{"tipi" = two dogs.} \\
\texttt{This implies "tipi" = "two dogs".} \\
\texttt{But "tipi" = "boat".} \\
\texttt{This suggests "tipi" is a number word "two"?} \\
\texttt{"two dogs" = "tipi" (two) + "dog" (stem)?} \\
\texttt{But where is the stem for dog?} \\
\texttt{Maybe "tipi" is "dog" and "dugu" is "boat" and the translation "a dog watch a boat" is "dugu deniru tipi" where "dugu" = dog and "tipi" = boat.} \\
\texttt{Then "two dogs watch a dog" => "tipi deniru tipinu nu".} \\
\texttt{"tipi" = two dogs.} \\
\texttt{This implies "tipi" = "two dogs".} \\
\texttt{But "tipi" = "boat".} \\
\texttt{This is a paradox.} \\
\texttt{Maybe the word for "dog" is "tipi" and "boat" is "dugu"?} \\
\texttt{"a dog watch a boat" => "dugu deniru tipi".} \\
\texttt{If "dugu" = dog, "tipi" = boat.} \\
\texttt{"two dogs watch a dog" => "tipi deniru tipinu nu".} \\
\texttt{"tipi" = two dogs.} \\
\texttt{This implies "tipi" = "two dogs".} \\
\texttt{But "tipi" = "boat".} \\
\texttt{Contradiction.} \\
\texttt{Maybe "tipi" is "dog" and "dugu" is "boat" and the translation "a dog watch a boat" is "dugu deniru tipi" where "dugu" = boat and "tipi" = dog?} \\
\texttt{Then "boat watch dog". Incorrect.} \\
\texttt{Maybe the word for "dog" is "dugu" and "boat" is "tipi" is correct.} \\
\texttt{Then "two dogs watch a dog" => "tipi deniru tipinu nu".} \\
\texttt{"tipi" = two dogs.} \\
\texttt{This implies "tipi" = "two dogs".} \\
\texttt{But "tipi" = "boat".} \\
\texttt{This suggests "tipi" is a homonym? Unlikely.} \\
\texttt{Maybe "tipi" is "dog" and "dugu" is "boat" and the translation "a dog watch a boat" is "dugu deniru tipi" where "dugu" = dog and "tipi" = boat.} \\
\texttt{Then "two dogs watch a dog" => "tipi deniru tipinu nu".} \\
\texttt{"tipi" = two dogs.} \\
\texttt{This implies "tipi" = "two dogs".} \\
\texttt{But "tipi" = "boat".} \\
\texttt{This is impossible.} \\
\texttt{Maybe the translation of "two dogs watch a dog" is "tipi deniru tipinu nu" where "tipi" = "two", "deniru" = "dogs watch", "tipinu" = "a dog", "nu" = marker?} \\
\texttt{No, "deniru" is "watch".} \\
\texttt{Maybe "tipi" is "dog" and "dugu" is "boat" and the translation "a dog watch a boat" is "dugu deniru tipi" where "dugu" = dog and "tipi" = boat.} \\
\texttt{Then "two dogs watch a dog" => "tipi deniru tipinu nu".} \\
\texttt{"tipi" = two dogs.} \\
\texttt{This implies "tipi" = "two dogs".} \\
\texttt{But "tipi" = "boat".} \\
\texttt{This is a contradiction.} \\
\texttt{Maybe the word for "dog" is "tipi" and "boat" is "dugu"?} \\
\texttt{"a dog watch a boat" => "dugu deniru tipi".} \\
\texttt{If "dugu" = dog, "tipi" = boat.} \\
\texttt{"two dogs watch a dog" => "tipi deniru tipinu nu".} \\
\texttt{"tipi" = two dogs.} \\
\texttt{This implies "tipi" = "two dogs".} \\
\texttt{But "tipi" = "boat".} \\
\texttt{Contradiction.} \\
\texttt{Maybe "tipi" is "dog" and "dugu" is "boat" and the translation "a dog watch a boat" is} \\
            }
        }
    }
    & & \\ \\

    \theutterance \stepcounter{utterance}  
    & & & \multicolumn{2}{p{0.3\linewidth}}{
        \cellcolor[rgb]{0.9,0.9,0.9}{
            \makecell[{{p{\linewidth}}}]{
                \texttt{\tiny{[GM$|$GM]}}
                \texttt{could not find any probe sentences; put them under a 'PROBES:' line, one per line} \\
            }
        }
    }
    & & \\ \\

\end{supertabular}
}

\end{document}
